% =====================================================================
%  STANDALONE. The bibliography is embedded below via filecontents, so
%  this single file is the whole document -- no external .bib needed.
%  Paste it into a fresh Overleaf project and compile.
%
%      pdflatex frequency_optimised_findings
%      bibtex   frequency_optimised_findings
%      pdflatex frequency_optimised_findings
%      pdflatex frequency_optimised_findings
%
%  or simply:  latexmk -pdf frequency_optimised_findings
%
%  SCOPE. Everything learned since the patch family moved from pixel-space
%  realism constraints to a FREQUENCY-DOMAIN perceptual constraint, written
%  as guidance for the remaining experiments. Companion to
%  docs/csf_patch_procedure.tex, which states the procedure formally; this
%  document states what the procedure MEASURED and what to do next.
%
%  PROVENANCE IS MARKED THROUGHOUT. Elements carrying a citation come from
%  the literature; those marked [no external source] are contributions of
%  this work and must be defended in the text rather than cited. Sec. 12
%  tabulates it, and Sec. 12.2 says which bibliography entries are verified
%  and which still need a check.
% =====================================================================
\documentclass[11pt,a4paper]{article}

\begin{filecontents*}{frequency_optimised_findings.bib}
% ---- perceptual model ------------------------------------------------
@inproceedings{barten2003,
  author    = {Barten, Peter G. J.},
  title     = {Formula for the contrast sensitivity of the human eye},
  booktitle = {Image Quality and System Performance, Proc. SPIE},
  volume    = {5294},
  pages     = {231--238},
  year      = {2003},
  doi       = {10.1117/12.537476}
}
@book{barten1999,
  author    = {Barten, Peter G. J.},
  title     = {Contrast Sensitivity of the Human Eye and Its Effects on
               Image Quality},
  publisher = {SPIE Press},
  address   = {Bellingham, WA},
  year      = {1999},
  doi       = {10.1117/3.353254}
}
@article{watson2005sso,
  author  = {Watson, Andrew B. and Ahumada, Albert J.},
  title   = {A standard model for foveal detection of spatial contrast},
  journal = {Journal of Vision},
  volume  = {5}, number = {9}, pages = {717--740}, year = {2005},
  doi     = {10.1167/5.9.6}
}
@article{watson2000five,
  author  = {Watson, Andrew B.},
  title   = {Visual detection of spatial contrast patterns: Evaluation of
             five simple models},
  journal = {Optics Express},
  volume  = {6}, number = {1}, pages = {12--33}, year = {2000},
  doi     = {10.1364/OE.6.000012}
}
@article{quick1974,
  author  = {Quick, R. F.},
  title   = {A vector-magnitude model of contrast detection},
  journal = {Kybernetik},
  volume  = {16}, number = {2}, pages = {65--67}, year = {1974},
  doi     = {10.1007/BF00271628}
}
@article{robson1981,
  author  = {Robson, J. G. and Graham, Norma},
  title   = {Probability summation and regional variation in contrast
             sensitivity across the visual field},
  journal = {Vision Research},
  volume  = {21}, number = {3}, pages = {409--418}, year = {1981},
  doi     = {10.1016/0042-6989(81)90169-3}
}
@article{watson1979,
  author  = {Watson, Andrew B.},
  title   = {Probability summation over time},
  journal = {Vision Research},
  volume  = {19}, number = {5}, pages = {515--522}, year = {1979},
  doi     = {10.1016/0042-6989(79)90136-6}
}
@book{michelson1927,
  author = {Michelson, Albert A.}, title = {Studies in Optics},
  publisher = {University of Chicago Press}, year = {1927}
}
@article{campbell1968,
  author  = {Campbell, F. W. and Robson, J. G.},
  title   = {Application of Fourier analysis to the visibility of gratings},
  journal = {The Journal of Physiology},
  volume  = {197}, number = {3}, pages = {551--566}, year = {1968},
  doi     = {10.1113/jphysiol.1968.sp008574}
}
@article{galinska2026csflow,
  author  = {Galinska, Malgorzata and Pogodzinski, Bart and Lenssen, Jan Eric},
  title   = {{CSFlow}: Aligning Flow Matching with Human Contrast Sensitivity},
  journal = {arXiv preprint arXiv:2606.08833},
  year    = {2026},
  note    = {identifier carried over from docs/csf\_patch\_procedure.tex}
}
@misc{iec61966,
  author = {{International Electrotechnical Commission}},
  title  = {{IEC 61966-2-1:1999} Multimedia systems and equipment ---
            Colour measurement and management --- Part 2-1: Colour
            management --- Default RGB colour space --- sRGB},
  year   = {1999}, howpublished = {International standard}
}
@techreport{itu709,
  author      = {{ITU-R}},
  title       = {Recommendation {BT.709-6}: Parameter values for the {HDTV}
                 standards for production and international programme exchange},
  institution = {International Telecommunication Union}, year = {2015}
}
% ---- image statistics ------------------------------------------------
@article{field1987,
  author  = {Field, David J.},
  title   = {Relations between the statistics of natural images and the
             response properties of cortical cells},
  journal = {Journal of the Optical Society of America A},
  volume  = {4}, number = {12}, pages = {2379--2394}, year = {1987},
  doi     = {10.1364/JOSAA.4.002379}
}
@article{falck2025fourier,
  author  = {Falck, Fabian and Pandeva, Teodora and Zahirnia, Kiarash and
             Lawrence, Rachel and Turner, Richard and Meeds, Edward and
             Zazo, Javier and Karmalkar, Sushrut},
  title   = {A Fourier Space Perspective on Diffusion Models},
  journal = {arXiv preprint arXiv:2505.11278}, year = {2025}
}
@misc{dieleman2024spectral,
  author = {Dieleman, Sander}, title = {Diffusion is spectral autoregression},
  year = {2024},
  howpublished = {Blog post, sander.ai/2024/09/02/spectral-autoregression.html}
}
@article{ruzanski2011rapsd,
  author  = {Ruzanski, Evan and Chandrasekar, V.},
  title   = {Scale Filtering for Improved Nowcasting Performance in a
             High-Resolution X-Band Radar Network},
  journal = {IEEE Transactions on Geoscience and Remote Sensing},
  volume  = {49}, number = {6}, pages = {2296--2307}, year = {2011},
  doi     = {10.1109/TGRS.2010.2103946}
}
% ---- attacks ---------------------------------------------------------
@inproceedings{goodfellow2015fgsm,
  author    = {Goodfellow, Ian J. and Shlens, Jonathon and Szegedy, Christian},
  title     = {Explaining and Harnessing Adversarial Examples},
  booktitle = {International Conference on Learning Representations (ICLR)},
  year      = {2015}
}
@inproceedings{madry2018pgd,
  author    = {Madry, Aleksander and Makelov, Aleksandar and Schmidt, Ludwig and
               Tsipras, Dimitris and Vladu, Adrian},
  title     = {Towards Deep Learning Models Resistant to Adversarial Attacks},
  booktitle = {International Conference on Learning Representations (ICLR)},
  year      = {2018}
}
@article{brown2017patch,
  author  = {Brown, Tom B. and Man{\'e}, Dandelion and Roy, Aurko and
             Abadi, Mart{\'\i}n and Gilmer, Justin},
  title   = {Adversarial Patch},
  journal = {arXiv preprint arXiv:1712.09665}, year = {2017}
}
@inproceedings{xie2017dag,
  author    = {Xie, Cihang and Wang, Jianyu and Zhang, Zhishuai and Zhou, Yuyin
               and Xie, Lingxi and Yuille, Alan},
  title     = {Adversarial Examples for Semantic Segmentation and Object
               Detection},
  booktitle = {IEEE International Conference on Computer Vision (ICCV)},
  year      = {2017}, doi = {10.1109/ICCV.2017.153}
}
@inproceedings{gu2022segpgd,
  author    = {Gu, Jindong and Zhao, Hengshuang and Tresp, Volker and
               Torr, Philip H. S.},
  title     = {{SegPGD}: An Effective and Efficient Adversarial Attack for
               Evaluating and Boosting Segmentation Robustness},
  booktitle = {European Conference on Computer Vision (ECCV)},
  year      = {2022}, doi = {10.1007/978-3-031-19818-2_18}
}
@inproceedings{agnihotri2024cospgd,
  author    = {Agnihotri, Shashank and Jung, Steffen and Keuper, Margret},
  title     = {{CosPGD}: An Efficient White-Box Adversarial Attack for
               Pixel-Wise Prediction Tasks},
  booktitle = {International Conference on Machine Learning (ICML)}, year = {2024}
}
@article{mirsky2023ipatch,
  author  = {Mirsky, Yisroel}, title = {{IPatch}: A Remote Adversarial Patch},
  journal = {Cybersecurity}, volume = {6}, number = {1}, year = {2023},
  doi     = {10.1186/s42400-023-00145-0}
}
@inproceedings{tan2021lap,
  author    = {Tan, Jia and Ji, Nan and Xie, Haidong and Xiang, Xueshuang},
  title     = {Legitimate Adversarial Patches: Evading Human Eyes and
               Detection Models in the Physical World},
  booktitle = {Proceedings of the 29th ACM International Conference on Multimedia},
  pages     = {5307--5315}, year = {2021}, doi = {10.1145/3474085.3475653}
}
@inproceedings{sharif2016nps,
  author    = {Sharif, Mahmood and Bhagavatula, Sruti and Bauer, Lujo and
               Reiter, Michael K.},
  title     = {Accessorize to a Crime: Real and Stealthy Attacks on
               State-of-the-Art Face Recognition},
  booktitle = {ACM SIGSAC Conference on Computer and Communications Security (CCS)},
  pages     = {1528--1540}, year = {2016}, doi = {10.1145/2976749.2978392}
}
@inproceedings{athalye2018eot,
  author    = {Athalye, Anish and Engstrom, Logan and Ilyas, Andrew and
               Kwok, Kevin},
  title     = {Synthesizing Robust Adversarial Examples},
  booktitle = {International Conference on Machine Learning (ICML)}, year = {2018}
}
@inproceedings{moosavi2017universal,
  author    = {Moosavi-Dezfooli, Seyed-Mohsen and Fawzi, Alhussein and
               Fawzi, Omar and Frossard, Pascal},
  title     = {Universal Adversarial Perturbations},
  booktitle = {IEEE Conference on Computer Vision and Pattern Recognition (CVPR)},
  year      = {2017}, doi = {10.1109/CVPR.2017.17}
}
@inproceedings{guo2020lowfreq,
  author    = {Guo, Chuan and Frank, Jared S. and Weinberger, Kilian Q.},
  title     = {Low Frequency Adversarial Perturbation},
  booktitle = {Uncertainty in Artificial Intelligence (UAI)}, year = {2020}
}
@inproceedings{sharma2019spectral,
  author    = {Sharma, Yash and Ding, Gavin Weiguang and Brubaker, Marcus A.},
  title     = {On the Effectiveness of Low Frequency Perturbations},
  booktitle = {International Joint Conference on Artificial Intelligence (IJCAI)},
  year      = {2019}, doi = {10.24963/ijcai.2019/470}
}
% ---- entropy ---------------------------------------------------------
@article{tsallis1988,
  author  = {Tsallis, Constantino},
  title   = {Possible generalization of {Boltzmann-Gibbs} statistics},
  journal = {Journal of Statistical Physics},
  volume  = {52}, number = {1--2}, pages = {479--487}, year = {1988},
  doi     = {10.1007/BF01016429}
}
@article{havrda1967,
  author  = {Havrda, Jan and Charv{\'a}t, Franti{\v s}ek},
  title   = {Quantification method of classification processes: Concept of
             structural $a$-entropy},
  journal = {Kybernetika}, volume = {3}, number = {1}, pages = {30--35},
  year    = {1967}
}
@article{shannon1948,
  author  = {Shannon, Claude E.},
  title   = {A Mathematical Theory of Communication},
  journal = {The Bell System Technical Journal},
  volume  = {27}, number = {3}, pages = {379--423}, year = {1948},
  doi     = {10.1002/j.1538-7305.1948.tb01338.x}
}
@inproceedings{grandvalet2004,
  author    = {Grandvalet, Yves and Bengio, Yoshua},
  title     = {Semi-supervised Learning by Entropy Minimization},
  booktitle = {Advances in Neural Information Processing Systems (NeurIPS)},
  year      = {2004}
}
@inproceedings{wang2021tent,
  author    = {Wang, Dequan and Shelhamer, Evan and Liu, Shaoteng and
               Olshausen, Bruno and Darrell, Trevor},
  title     = {Tent: Fully Test-Time Adaptation by Entropy Minimization},
  booktitle = {International Conference on Learning Representations (ICLR)},
  year      = {2021}
}
% ---- models, data, optimisation -------------------------------------
@inproceedings{xie2021segformer,
  author    = {Xie, Enze and Wang, Wenhai and Yu, Zhiding and Anandkumar, Anima
               and Alvarez, Jose M. and Luo, Ping},
  title     = {{SegFormer}: Simple and Efficient Design for Semantic
               Segmentation with Transformers},
  booktitle = {Advances in Neural Information Processing Systems (NeurIPS)},
  year      = {2021}
}
@inproceedings{cordts2016cityscapes,
  author    = {Cordts, Marius and Omran, Mohamed and Ramos, Sebastian and
               Rehfeld, Timo and Enzweiler, Markus and Benenson, Rodrigo and
               Franke, Uwe and Roth, Stefan and Schiele, Bernt},
  title     = {The Cityscapes Dataset for Semantic Urban Scene Understanding},
  booktitle = {IEEE Conference on Computer Vision and Pattern Recognition (CVPR)},
  year      = {2016}, doi = {10.1109/CVPR.2016.350}
}
@inproceedings{kingma2015adam,
  author    = {Kingma, Diederik P. and Ba, Jimmy},
  title     = {Adam: A Method for Stochastic Optimization},
  booktitle = {International Conference on Learning Representations (ICLR)},
  year      = {2015}
}
@inproceedings{loshchilov2017sgdr,
  author    = {Loshchilov, Ilya and Hutter, Frank},
  title     = {{SGDR}: Stochastic Gradient Descent with Warm Restarts},
  booktitle = {International Conference on Learning Representations (ICLR)},
  year      = {2017}
}
@inproceedings{luo2016erf,
  author    = {Luo, Wenjie and Li, Yujia and Urtasun, Raquel and Zemel, Richard},
  title     = {Understanding the Effective Receptive Field in Deep
               Convolutional Neural Networks},
  booktitle = {Advances in Neural Information Processing Systems (NeurIPS)},
  year      = {2016}
}
@inproceedings{zhang2018lpips,
  author    = {Zhang, Richard and Isola, Phillip and Efros, Alexei A. and
               Shechtman, Eli and Wang, Oliver},
  title     = {The Unreasonable Effectiveness of Deep Features as a
               Perceptual Metric},
  booktitle = {IEEE Conference on Computer Vision and Pattern Recognition (CVPR)},
  year      = {2018}
}
@book{efron1993bootstrap,
  author    = {Efron, Bradley and Tibshirani, Robert J.},
  title     = {An Introduction to the Bootstrap},
  publisher = {Chapman and Hall/CRC}, year = {1993},
  doi       = {10.1201/9780429246593}
}
\end{filecontents*}

% lmodern before fontenc: with T1 and no scalable font, microtype's font
% expansion aborts with "auto expansion is only possible with scalable
% fonts" and no PDF is produced.
\usepackage{lmodern}
\usepackage[T1]{fontenc}
\usepackage[utf8]{inputenc}
\usepackage[a4paper,margin=2.5cm]{geometry}
\usepackage{amsmath,amssymb}
\usepackage{booktabs}
\usepackage{array}
\usepackage{longtable}
\usepackage{microtype}
\usepackage[round,authoryear]{natbib}
\usepackage[hidelinks]{hyperref}

\newcommand{\novel}{\textbf{[no external source]}}
\newcolumntype{L}[1]{>{\raggedright\arraybackslash}p{#1}}
\newcommand{\cpp}{\mathrm{cyc/px}}
\newcommand{\cpd}{\mathrm{cpd}}
\newcommand{\sg}[1]{\left[#1\right]_{\mathrm{sg}}}

\title{Frequency-Optimised Adversarial Patches for Semantic
       Segmentation:\\[2pt]
       \large What Has Been Measured, and What to Run Next}
\author{Findings, formal statement, and experimental design}
\date{}

\begin{document}
\maketitle

\begin{abstract}
\noindent
This document collects what has been established since the patch family
moved from pixel-space realism constraints to a constraint stated in the
\emph{frequency domain} by a model of human contrast sensitivity, and turns
it into guidance for the remaining experiments. Three results reframe the
premise. First, the segmentation network is approximately \emph{spectrally
flat}: it reads every radial band up to Nyquist about equally, so the
exploitable asymmetry is not a band that the network reads and the eye
cannot see, but a $\sim$47$\times$ difference in what a unit of prediction
change \emph{costs} perceptually. Second, the perceptual budget in use was
mis-calibrated by a factor of $\approx1.9$ at the frequencies that matter,
for two independent and partly cancelling reasons that are derived in closed
form in Sec.~4.4. Third, a constraint that was believed to hold by
construction was in fact violated by $2.8$--$7.4\times$ during optimisation,
because clipping against real image content is a hard nonlinearity that
generates broadband harmonics; enforcing it costs $4.6$ mIoU points and is
now mandatory for any number quoted with a $\tau$ attached. Section~8 states
the three untargeted objectives --- cross-entropy, CosPGD, and a
Tsallis-entropy objective --- with full derivations, including the closed
form of the Tsallis gradient and its Shannon limit. Every symbol in every
equation is defined where it first appears. Elements without an external
source are marked \novel{} and tabulated in Sec.~12.
\end{abstract}

\tableofcontents

% =====================================================================
\section{Scope: what changed, and why this document exists}

The patch family in this project has three generations. The first
parameterised a patch in \emph{pixel space} and constrained nothing
\citep{brown2017patch}. The second added \emph{pixel-space realism}: a pull
toward a reference image, a total-variation penalty, and printability, after
\citet{tan2021lap} and \citet{sharif2016nps}. Neither generation knows
anything about an observer --- both penalise statistics of the patch rather
than its \emph{visibility}.

The third generation, which this document is about, states the constraint in
the frequency domain against a contrast sensitivity function (CSF). That
converts ``make the patch look plausible'' into a falsifiable statement: a
perturbation is admissible if its contrast at every spatial frequency sits
below the detection threshold of a modelled observer at that frequency
\citep{campbell1968,barten2003}.

The motivating asymmetry was stated as a hypothesis and has now been
measured. It was: a segmentation network has no CSF and weights the band up
to Nyquist uniformly, whereas the human visual system does not, so a
perturbation may be admissible to one and not to the other. The named risk
was that SegFormer's patch embedding has stride~4 \citep{xie2021segformer},
so the highest frequencies might be attenuated before the first attention
block ever sees them --- in which case the attack would be bounded by an
architectural detail rather than by the perceptual budget.

\paragraph{The premise survives, but not in its original form.}
Section~10.1 gives the measurement. The network reads \emph{everything},
including the high band; what differs across bands is not effectiveness but
perceptual price. The claim to make in the thesis is therefore about
\emph{efficiency}, and never about a band the network can see and the eye
cannot.

\paragraph{Relation to the companion document.}
\texttt{docs/csf\_patch\_procedure.tex} states the procedure formally and
carries its own provenance table. This document does not repeat those
derivations except where a measurement has changed them; where the two
disagree, this one is later.

\paragraph{How to read the findings.}
Findings are labelled \textbf{F1}--\textbf{F9} in Sec.~10 and are referred to
by label from the design guidance in Sec.~11. Each states the measurement,
the mechanism, and --- most importantly --- the claim it does \emph{not}
support.

% =====================================================================
\section{Notation}

Every symbol used anywhere below, in one place.

\begin{center}
\small
\begin{longtable}{@{}lL{10.9cm}@{}}
\toprule
\multicolumn{2}{@{}l}{\textbf{Images and the network}}\\
\midrule
$\mathbf{x}\in\mathbb{R}^{3\times H\times W}$ & input image, ImageNet-normalised, as the network receives it \\
$\mathbf{x}^{01}$ & the same image in $[0,1]$ sRGB \emph{code} values, $\mathrm{clip}(\mathbf{x}\odot\boldsymbol{\sigma}_n+\boldsymbol{\mu}_n,0,1)$, where $\boldsymbol{\mu}_n,\boldsymbol{\sigma}_n$ are the per-channel ImageNet statistics \\
$H,W$ & image height and width in pixels; $512\times1024$ throughout unless stated \\
$\mathbf{y}$ & ground-truth label map with trainIds in $\{0,\dots,C-1\}\cup\{255\}$, where $255$ is \emph{void} and is never scored \\
$C$ & number of classes; $C=19$ for Cityscapes \citep{cordts2016cityscapes} \\
$f$ & the frozen segmentation network; $f(\mathbf{x})\in\mathbb{R}^{C\times H'\times W'}$ are logits at the head's own resolution \\
$\mathcal{U}$ & bilinear upsampling from $H'\times W'$ to $H\times W$ \\
$\mathbf{p}_{uv}$ & $=\mathrm{softmax}\big(\mathcal{U}f(\cdot)_{:,u,v}\big)$, the predicted class distribution at pixel $(u,v)$, a point in the simplex $\Delta^{C-1}$ \\
$\mathbf{e}_k$ & the one-hot vector of class $k$ in $\mathbb{R}^C$ \\
\midrule
\multicolumn{2}{@{}l}{\textbf{The patch}}\\
\midrule
$s$ & patch scale as a fraction of image height; $s=0.25$ \\
$p=\lfloor Hs\rfloor$ & rendered patch side in pixels; $p=128$ at $H=512$, $s=0.25$ \\
$S$ & parameter-grid resolution (\texttt{--patch\_size}). \textbf{Must equal $p$} in frequency modes (Sec.~6.4) \\
$\mathbf{z}\in\mathbb{R}^{3\times S\times S}$ & the optimised parameter, unconstrained \\
$\boldsymbol{\delta}\in\mathbb{R}^{3\times S\times S}$ & the residual: a \emph{signed} additive perturbation in sRGB code units \\
$\mathbf{r}$ & the residual's \emph{base}: the image content the patch covers (\texttt{--from\_image}), or mid grey \\
$\mathcal{M}$ & silhouette mask, $1$ on pixels actually composited; all ones for a square patch \\
$\mathbf{m}$ & the footprint: $\mathcal{M}$ padded into the full frame \\
\midrule
\multicolumn{2}{@{}l}{\textbf{Frequency}}\\
\midrule
$\mathcal{F},\mathcal{F}^{-1}$ & real 2-D DFT (\texttt{rfft2}) and its inverse, \emph{forward-normalised}: $\mathcal{F}$ carries the $1/(S^2)$ factor, so a cosine of peak amplitude $A$ has coefficient magnitude $A/2$ \\
$\mathcal{K}$ & the half-spectrum bin index set of an $S\times S$ real signal; $|\mathcal{K}|=S(\lfloor S/2\rfloor+1)$ \\
$f_y,f_x$ & bin frequencies in cycles per pixel, each in $[-\tfrac12,\tfrac12]$ \\
$\rho=\sqrt{f_y^2+f_x^2}$ & radial frequency in cycles per pixel ($\cpp$); $\rho_{\max}=1/2$ is Nyquist \\
$\theta$ & degrees of visual angle subtended by one pixel (Sec.~3.1) \\
$u=\rho/\theta$ & radial frequency in cycles per degree ($\cpd$) --- the unit a CSF is defined in \\
\midrule
\multicolumn{2}{@{}l}{\textbf{The perceptual constraint}}\\
\midrule
$\mathrm{CSF}(u)$ & contrast sensitivity: the reciprocal of the just-detectable Michelson contrast at $u$ \\
$\kappa$ & contrast scale: converts a DFT magnitude into a Michelson contrast (Sec.~4.1) \\
$B(f)$ & per-bin amplitude budget (Sec.~4.2) \\
$A_{\max}$ & hard cap on $B$ where $\mathrm{CSF}\to0$; $A_{\max}=0.25$ \\
$\beta$ & Minkowski exponent modelling probability summation; $\beta=3$ \\
$V_\beta(\boldsymbol{\delta})$ & visibility index of $\boldsymbol{\delta}$, in nominal JND (Sec.~5.2) \\
$\tau$ & the perceptual budget: the permitted value of $V_\beta$. \textbf{The single knob of this attack family} \\
$Y$ & linear relative luminance in $[0,1]$; $Y_{\mathrm{ref}}$ is the reference the CSF is evaluated at \\
$c$ & an sRGB code value in $[0,1]$; $c_{\mathrm{ref}}$ is the code value whose luminance is $Y_{\mathrm{ref}}$ \\
$L$ & absolute luminance in $\mathrm{cd/m^2}$; $L=Y\,L_{\mathrm{peak}}$ with $L_{\mathrm{peak}}=100$ \\
$E$ & retinal illuminance in trolands (Td) \\
\midrule
\multicolumn{2}{@{}l}{\textbf{Evaluation}}\\
\midrule
$\Omega$ & the scored support: valid, non-footprint, optionally reach-restricted pixels (Sec.~8.1) \\
$\Delta\mathrm{mIoU}_{\mathrm{rem}}$ & remote mIoU drop, the headline attack metric (Sec.~9.1) \\
$\phi$ & any-flip rate: fraction of remote pixels whose $\arg\max$ changed (Sec.~9.2) \\
$\eta$ & spectral efficiency: prediction change bought per unit of visibility (Sec.~9.4) \\
\bottomrule
\end{longtable}
\end{center}

% =====================================================================
\section{The observer model}

A perceptual constraint is only as meaningful as the observer it assumes.
This section fixes that observer completely: the geometry, the sensitivity
function, and the transfer function relating the numbers the optimiser
touches to the light a human would receive.

\subsection{Viewing geometry: the assumption everything is conditional on}

A CSF is defined over cycles per \emph{degree of visual angle}. An image
spectrum is in cycles per \emph{pixel}. Converting between them requires a
physical pixel size and a viewing distance:
\begin{align}
  \theta &= \frac{180}{\pi}\cdot 2\arctan\!\left(\frac{P}{2D}\right)
     \qquad\text{[degrees per pixel]}, \label{eq:theta}\\
  u &= \rho/\theta \qquad\text{[cycles per degree]}. \label{eq:cpd}
\end{align}
Here $P$ is the physical size of one displayed pixel in cm, $D$ is the
distance from observer to display in cm, $\rho$ is radial frequency in
$\cpp$ from Sec.~2, and $u$ is the frequency the CSF is evaluated at. The
factor $180/\pi$ converts radians to degrees; the factor $2$ with the
half-angle $P/2D$ is the standard small-object subtense.

\paragraph{Only the ratio $P/D$ matters.}
Equation~\eqref{eq:theta} depends on $P$ and $D$ only through $P/(2D)$, so
they are not two independent knobs but one. A geometry can therefore be
quoted as a single number, $\theta$ in deg/px, and any $(P,D)$ pair
realising it is equivalent. Report $\theta$; never report $P$ and $D$ as
though they were separate degrees of freedom.

\paragraph{Defaults, and why they are not the threat model.}
The defaults are $P=0.0114$~cm and $D=50$~cm, reproducing the geometry of
\citet[App.~D]{galinska2026csflow}, i.e.\ a laptop screen. They give
$\theta\approx0.01306$ deg/px and a Nyquist frequency of
\begin{equation}
  u_{\mathrm{Nyq}} = \frac{0.5}{\theta} \approx 38.3\ \cpd ,
  \label{eq:nyquist}
\end{equation}
which is \emph{below} the $\approx50$--$60$ $\cpd$ where human sensitivity
actually vanishes. Consequence, and it must be stated: the sampling grid
cannot represent a truly invisible frequency. The budget at Nyquist is large
but finite, and ``invisible'' here always means ``below threshold'', never
``outside the passband''.

\paragraph{Two constructors, because deriving $P$ by hand is where the
mistake happens.} For a real display of diagonal $d_{\mathrm{in}}$ inches,
$N$ pixels across, and aspect ratio $a{:}b$,
\begin{equation}
  P = \frac{2.54\, d_{\mathrm{in}}\, a}{N\sqrt{a^2+b^2}} .
  \label{eq:fromscreen}
\end{equation}
A 15.6-inch 1080p laptop gives $P=0.0180$~cm, which is $1.6\times$ the
default; a residual built for the default geometry and inspected on that
laptop at 50~cm measures roughly $5\times$ its requested visibility. Nothing
in the method is wrong there --- the geometry is simply describing a
different monitor.

For a \emph{physical} patch of width $w$ metres, viewed at distance $d$
metres, rendered on an $S\times S$ grid,
\begin{equation}
  \theta = \frac{1}{S}\cdot\frac{180}{\pi}\cdot
           2\arctan\!\left(\frac{w}{2d}\right).
  \label{eq:fromphysical}
\end{equation}
A $0.5$~m patch at $20$~m on a $128$~px grid gives $\theta=0.01119$ deg/px,
i.e.\ $u_{\mathrm{Nyq}}=44.7$ $\cpd$: a \emph{finer} geometry than the desk
monitor. The road threat model therefore has \emph{more} room to hide in
than inspecting PNGs on a laptop suggests, which is the opposite of the
intuition one gets from looking at saved images.

\paragraph{Validation by eye is a consistency check, not a perceptual
study.} A patch built for geometry $\theta$ may be inspected on any display
of pixel density $\mathrm{ppi}$ provided it is viewed at
\begin{equation}
  D_{\mathrm{match}} = \frac{2.54/\mathrm{ppi}}
                            {2\tan(\theta_{\mathrm{rad}}/2)}
  \;\propto\; \frac{1}{\mathrm{ppi}},
  \label{eq:matched}
\end{equation}
with $\theta_{\mathrm{rad}}$ the per-pixel angle in radians. At the default
geometry that is $79$~cm on a 141~PPI laptop and $119$~cm on a 93~PPI
desktop panel. Viewing closer is not a stricter test of the same claim; it
is a test of a different claim.

\subsection{Barten's CSF}

The default sensitivity model is Barten's \citeyearpar{barten2003}
formula, with the parameter values of \citet[App.~A]{galinska2026csflow}:
\begin{equation}
  \mathrm{CSF}(u) =
  \frac{M_{\mathrm{opt}}(u)/k}
       {\sqrt{\dfrac{2}{T}\left(\dfrac{1}{X_0^2}+\dfrac{1}{X_{\max}^2}
              +\dfrac{u^2}{N_{\max}^2}\right)
              \left(\dfrac{1}{\eta_q p_{\!\gamma} E}
              +\dfrac{\Phi_0}{1-e^{-(u/u_0)^2}}\right)}} .
  \label{eq:barten}
\end{equation}
\begin{center}
\small
\begin{tabular}{@{}llL{7.7cm}@{}}
\toprule
symbol & value & meaning \\
\midrule
$M_{\mathrm{opt}}(u)$ & --- & optical modulation transfer of the eye, Eq.~\eqref{eq:mtf} \\
$k$ & $3.0$ & signal-to-noise ratio the visual system requires for detection \\
$T$ & $0.1$ s & integration time of the eye \\
$X_0$ & $60^\circ$ & angular size of the object \\
$X_{\max}$ & $12^\circ$ & maximum angular area the eye integrates over \\
$N_{\max}$ & $15$ & maximum number of cycles integrated \\
$\eta_q$ & $0.03$ & quantum efficiency of the eye \\
$p_{\!\gamma}$ & $1.2\times10^{6}$ & photon conversion factor, photons\,s$^{-1}$deg$^{-2}$Td$^{-1}$ \\
$E$ & $500$ Td & retinal illuminance (\textbf{revised in Sec.~4.4}) \\
$\Phi_0$ & $3\times10^{-8}$ deg$^2$s & spectral density of neural noise \\
$u_0$ & $7$ $\cpd$ & lateral-inhibition cutoff \\
$\sigma$ & $0.5/60$ deg & standard deviation of the optical point spread \\
\bottomrule
\end{tabular}
\end{center}
The optical transfer factor is
\begin{equation}
  M_{\mathrm{opt}}(u)=\exp\!\left(-2\pi^2\sigma^2u^2\right).
  \label{eq:mtf}
\end{equation}

\paragraph{Reading the formula.}
The three factors have distinct roles. $M_{\mathrm{opt}}$ is the eye's
optics and is a Gaussian low-pass: it is what kills sensitivity at high
$u$. The first bracket under the root is \emph{spatio-temporal integration}
--- how much signal the visual system can accumulate --- and grows like
$u^2$, which is a second high-frequency penalty. The second bracket is
\emph{noise}: a photon-noise term $1/(\eta_q p_{\!\gamma}E)$ that shrinks as
the scene gets brighter, plus a neural-noise term divided by a
lateral-inhibition factor $1-e^{-(u/u_0)^2}$ that vanishes at DC.

\paragraph{$\mathrm{CSF}(0)=0$, and that is correct.}
As $u\to0$ the inhibition factor $1-e^{-(u/u_0)^2}\to0$, the neural-noise
term diverges, and sensitivity goes to zero. Physically this says a uniform
field has no contrast to detect. It is handled explicitly in code rather
than being left to produce a NaN, because $1/\mathrm{CSF}$ is what the
budget of Sec.~4.2 is built from and a NaN there poisons everything
downstream silently.

\paragraph{Sensitivity is a reciprocal threshold.}
$\mathrm{CSF}(u)=700$ means a grating of Michelson contrast $1/700$ at
frequency $u$ is just detectable. This is the definition that makes
Sec.~4.2 an inversion rather than an analogy.

\subsection{The Standard Spatial Observer, as a control}

The alternative model is the Standard Spatial Observer
\citep{watson2005sso,watson2000five}, which unlike Barten is
\emph{orientation-dependent}:
\begin{align}
  \mathrm{CSF}_{\mathrm{SSO}}(f_x,f_y) &=
      \mathrm{CSF}_{\mathrm{T}}(u)\cdot \mathrm{OE}(f_x,f_y),
      \label{eq:sso}\\
  \mathrm{CSF}_{\mathrm{T}}(u) &= g\left(e^{-u/f_m}
      - \ell\, e^{-u^2/s_o^2}\right), \label{eq:tyler}\\
  \mathrm{OE}(f_x,f_y) &= 1 - w\,
      \frac{4\left(1-e^{-u/o_s}\right)f_x^2f_y^2}{u^4}. \label{eq:oblique}
\end{align}
Here $u=\sqrt{f_x^2+f_y^2}$ in $\cpd$, $g=330.74$ is an overall gain,
$f_m=7.28$ sets the exponential high-frequency decay, $\ell=0.837$ is the
low-frequency loss and $s_o=1.809$ the scale at which that loss is
attenuated, $w=1$ weights the oblique effect ($w=0$ disables it) and
$o_s=6.664$ is its scale.

\paragraph{The oblique effect is a free asymmetry.}
$\mathrm{OE}=1$ on the axes ($f_x=0$ or $f_y=0$) and dips to a minimum at
$45^\circ$, where $f_x^2f_y^2/u^4=1/4$. Diagonal gratings are harder to see
than axis-aligned ones of identical amplitude, so a diagonal carrier is
cheaper perceptually at equal amplitude. This is not exploited by the
current attack and is a legitimate future rung.

\paragraph{Why keep two models.}
So that the visibility claim can be shown not to depend on one particular
CSF. If Barten and SSO disagree about a perturbation, that disagreement
belongs in the write-up. SSO has a second use as a control: it has \emph{no
luminance parameter} (it fits a single photopic observer), so under SSO only
the contrast denominator responds to background luminance and the two
luminance effects of Sec.~4.4 are cleanly separated.

\subsection{sRGB, luminance, and the unit bridge}

The optimiser perturbs sRGB \emph{code values}; contrast is defined on
\emph{linear luminance}. Conflating the two is the single easiest way to be
wrong by a power of $2.4$, and Sec.~4.4 shows it was worth a factor of
$1.9$ in the budget.

The sRGB transfer function and its inverse \citep{iec61966} are
\begin{align}
  Y(c) &= \begin{cases}
      c/12.92, & c\le0.04045,\\[2pt]
      \left(\dfrac{c+0.055}{1.055}\right)^{2.4}, & c>0.04045,
  \end{cases} \label{eq:srgb2lin}\\[4pt]
  c(Y) &= \begin{cases}
      12.92\,Y, & Y\le0.0031308,\\[2pt]
      1.055\,Y^{1/2.4}-0.055, & Y>0.0031308,
  \end{cases} \label{eq:lin2srgb}
\end{align}
with $c$ the code value in $[0,1]$ and $Y$ the linear value in $[0,1]$. The
linear segment below $0.04045$ is not decoration: a naive $c^{2.4}$ departs
from the standard by tens of percent in exactly the dark codes that
Cityscapes asphalt occupies.

The \emph{local gain} of that curve is the unit bridge:
\begin{equation}
  \frac{dY}{dc}\bigg|_{c} = \begin{cases}
      1/12.92, & c\le0.04045,\\[4pt]
      \dfrac{2.4}{1.055}\left(\dfrac{c+0.055}{1.055}\right)^{1.4},
        & c>0.04045.
  \end{cases}
  \label{eq:slope}
\end{equation}
A budget derived in linear units is not a bound on the tensor the optimiser
touches until it has been divided by this slope. The linearisation is
first-order and therefore valid for small perturbations; at $\tau$ of order
one the residual is far below one code step and this is comfortable, at
large $\tau$ it is not, and there the \emph{measured} realised visibility
rather than the budget is the number to trust.

Relative luminance of a colour pixel uses the Rec.~709 weights
\citep{itu709}, applied to \emph{linearised} channels:
\begin{equation}
  Y = 0.2126\,Y(c_R) + 0.7152\,Y(c_G) + 0.0722\,Y(c_B).
  \label{eq:luma}
\end{equation}
Linearise \emph{then} weight, never the reverse: averaging gamma-encoded
channels over-states the luminance of a dark region, which would flatter
every fixed-reference simplification in Sec.~10.3.

Finally, background luminance enters Barten's CSF through the pupil
\citep[Eq.~2.9]{barten1999}:
\begin{align}
  d_{\mathrm{pupil}}(L) &= 5 - 3\tanh\!\left(0.4\log_{10}L\right)
     \quad\text{[mm]}, \label{eq:pupil}\\
  E(L) &= \frac{\pi\, d_{\mathrm{pupil}}(L)^2}{4}\, L \quad\text{[Td]},
     \label{eq:trolands}
\end{align}
with $L$ the absolute background luminance in $\mathrm{cd/m^2}$. The pupil
ranges from $\approx8$~mm in near-darkness to $\approx2$~mm in bright light.
Equation~\eqref{eq:trolands} is the retinal illuminance $E$ that appears in
Eq.~\eqref{eq:barten}. This is the \emph{entire} mechanism by which
background luminance changes Barten's sensitivity, and it is why that
dependence is comparatively flat: the pupil partly compensates.

% =====================================================================
\section{The budget: inverting a sensitivity into an amplitude bound}

\subsection{From a DFT magnitude to a Michelson contrast}
\label{sec:kappa}

A DFT magnitude is not a contrast, and the gap is a constant that must be
written down rather than absorbed silently. A real cosine of peak amplitude
$A$ riding on a mean level $\mu$ has Michelson contrast \citep{michelson1927}
\begin{equation}
  c_{\mathrm{M}}
  = \frac{I_{\max}-I_{\min}}{I_{\max}+I_{\min}}
  = \frac{(\mu+A)-(\mu-A)}{(\mu+A)+(\mu-A)}
  = \frac{A}{\mu},
  \label{eq:michelson}
\end{equation}
where $I_{\max},I_{\min}$ are the maximum and minimum levels of the grating.
With the forward-normalised transform of Sec.~2 that cosine has coefficient
magnitude $|\mathcal{F}\boldsymbol{\delta}(f)| = A/2$. Substituting,
\begin{equation}
  c_{\mathrm{M}}(f) = \kappa\,\big|\mathcal{F}\boldsymbol{\delta}(f)\big|,
  \qquad
  \kappa = \frac{2}{\mu}.
  \label{eq:kappa}
\end{equation}
$\kappa$ is the \emph{contrast scale}. The legacy convention takes
$\mu=\tfrac12$ (mid grey), giving
\begin{equation}
  \kappa_{\mathrm{legacy}} = 4 .
  \label{eq:kappa4}
\end{equation}

\paragraph{This convention is wrong on this dataset, and Sec.~4.4 says by how
much.} It is retained as the default \emph{deliberately}, so that the
existing $\tau$ ladder keeps its meaning and every past run is restatable by
a constant factor rather than invalidated. Every results table must carry
both a nominal and a calibrated column.

\subsection{The amplitude budget}

A component of contrast $c$ at frequency $u$ is at the detection threshold
when $c\cdot\mathrm{CSF}(u)=1$, by the definition of sensitivity in Sec.~3.2.
Requiring instead that each component reach at most $\tau$ just-noticeable
differences and solving Eq.~\eqref{eq:kappa} for the magnitude gives the
per-bin \emph{amplitude budget} \novel{}
\begin{equation}
  B(f) \;=\; \min\!\left(
     \frac{\tau}{\kappa\,\mathrm{CSF}\big(u(f)\big)},\;
     A_{\max}\right),
  \label{eq:budget}
\end{equation}
where $f\in\mathcal{K}$ indexes a half-spectrum bin, $u(f)$ is that bin's
radial frequency in $\cpd$ via Eq.~\eqref{eq:cpd}, $\tau$ is the perceptual
budget, $\kappa$ is from Eq.~\eqref{eq:kappa}, and $A_{\max}=0.25$ caps the
budget where $\mathrm{CSF}\to0$ and $1/\mathrm{CSF}$ would diverge. A quarter
of the dynamic range in a single Fourier component is already far beyond
anything this constraint is meant to permit, so the cap binds on no bin
until $\tau\gtrsim4$.

\paragraph{The shape of $B$ is the whole premise.}
$B$ is proportional to $1/\mathrm{CSF}$, so it is small where the eye is
sensitive (a few $\cpd$) and large near Nyquist. That inversion is what
``camouflage in the frequency domain'' means operationally. The sanity check
before trusting any run: if $B$ at high frequency is not orders of magnitude
larger than $B$ at the CSF peak, there is no gap to hide in and the premise
fails at step zero.

\subsection{The low-frequency cutoff: making a full-field CSF apply to a
patch}
\label{sec:mincycles}

Equation~\eqref{eq:budget} as written has a defect that inverts its
intent. Since $\mathrm{CSF}(0)=0$, the bins nearest DC receive the
\emph{largest} allowance of any frequency. Measured: a residual built that
way put $3\times10^{-5}$ of its power in the lowest radial bin and
$10^{-9}$ in the highest, and the diagnostic table printed a budget of
$0.0770$ for $0.00$--$0.02\ \cpp$ against $0.0026$ at Nyquist --- the exact
opposite of what the attack is supposed to do.

This is not a bug in the CSF; it is a category error in applying it. The
low-frequency roll-off is measured for \emph{full-field} gratings, where the
eye adapts to a slow luminance ramp. Inside a patch of side $S$ there is no
such grating: a component below $1/S\ \cpp$ does not complete a single cycle
across the patch, so it is not a grating but a brightness offset --- and a
patch whose mean differs from its surroundings is a visible rectangle
whatever the CSF says about DC, because the patch border turns any offset
into an edge, and edges are broadband.

The fix is a hard low-frequency cutoff \novel{}:
\begin{equation}
  B(f) \;\leftarrow\; B(f)\cdot
     \mathbb{1}\!\left[\rho(f) \ge \frac{n_c}{S}\right],
  \qquad n_c = 2 ,
  \label{eq:mincycles}
\end{equation}
where $\rho(f)$ is the bin's radial frequency in $\cpp$, $S$ is the patch
side in pixels, $n_c$ is the minimum number of full cycles a component must
complete across the patch, and $\mathbb{1}[\cdot]$ is the indicator
function. Setting $n_c=0$ restores the naive behaviour and is useful only
for reproducing the failure. This cutoff carries over unchanged to every
budget variant below; nothing about luminance or universality changes the
argument for it.

\subsection{Calibration: the budget was $\approx1.9\times$ too permissive}
\label{sec:calib}

Two of the constants above are wrong on Cityscapes, and they are wrong in
opposite directions, so quoting only the net error hides the mechanism.

\paragraph{Error 1: the Michelson denominator.}
$\kappa=4$ assumes $\mu=0.5$. Measured over the full Cityscapes train split
($n=2975$) with a 128\,px footprint at $512\times1024$, the median mean code
value inside the footprint is $0.303$ at centre placement and $0.292$ on the
road surface at $(0.75,0.5)$; the val split gives $0.311$. In linear
relative luminance the medians are $Y=0.0971$, $0.0792$ and $0.1029$. So
$\mu=0.5$ is $1.65\times$ too bright at centre and $1.71\times$ on the road.
Contrast is $A/\mu$, so assuming a brighter background than the truth
\emph{under-states} contrast and over-states the admissible amplitude.

\paragraph{Error 2: the code-value/luminance conflation.}
$\kappa=2/\mu$ implicitly treats code values as luminances. They are not.
Redoing Eq.~\eqref{eq:kappa} correctly: a cosine of peak amplitude
$A_{\mathrm{code}}$ in \emph{code} units produces a luminance modulation
$A_{\mathrm{lin}} = (dY/dc)\,A_{\mathrm{code}}$ to first order, and
Michelson contrast is $A_{\mathrm{lin}}/Y_{\mathrm{ref}}$. Hence the
\emph{calibrated} contrast scale
\begin{equation}
  \kappa_{\mathrm{cal}}(Y_{\mathrm{ref}})
  = \frac{2}{Y_{\mathrm{ref}}}\cdot
    \frac{dY}{dc}\bigg|_{c_{\mathrm{ref}}},
  \qquad c_{\mathrm{ref}} = c(Y_{\mathrm{ref}}),
  \label{eq:kappacal}
\end{equation}
with $c(\cdot)$ from Eq.~\eqref{eq:lin2srgb} and the slope from
Eq.~\eqref{eq:slope}. \textbf{The legacy constant is not a special case of
this one.} At $Y_{\mathrm{ref}}=0.5$, Eq.~\eqref{eq:kappacal} returns
$\approx6.07$, not $4$, because a mid-\emph{luminance} grey is code $0.735$
where the transfer curve is steep.

\paragraph{Error 3, pushing the other way: retinal illuminance.}
$E=500$~Td in Eq.~\eqref{eq:barten} models a brighter-adapted and therefore
\emph{more sensitive} observer than the real $\approx115$~Td implied by
$Y=0.097$ at a $100\ \mathrm{cd/m^2}$ peak white through
Eqs.~\eqref{eq:pupil}--\eqref{eq:trolands}. Correcting it \emph{reduces}
$\mathrm{CSF}$ and therefore \emph{raises} the admissible amplitude, partly
cancelling errors 1 and 2.

\paragraph{The corrected budget, in closed form.}
Substituting Eq.~\eqref{eq:kappacal} into Eq.~\eqref{eq:budget} and
converting to code units by dividing by the slope gives a budget in linear
units and then in sRGB units:
\begin{align}
  B_{\mathrm{lin}}(f)
    &= \frac{\tau\,Y_{\mathrm{ref}}}
            {2\,\mathrm{CSF}\big(u(f);Y_{\mathrm{ref}}\big)},
    \label{eq:budgetlin}\\[3pt]
  B_{\mathrm{sRGB}}(f)
    &= \frac{B_{\mathrm{lin}}(f)}{(dY/dc)|_{c_{\mathrm{ref}}}}
     = \frac{\tau\left(c_{\mathrm{ref}}+0.055\right)}
            {4.8\;\mathrm{CSF}\big(u(f);Y_{\mathrm{ref}}\big)} .
  \label{eq:budgetsrgb}
\end{align}
The second equality is exact for $c_{\mathrm{ref}}>0.04045$: writing
$v=(c_{\mathrm{ref}}+0.055)/1.055$, Eq.~\eqref{eq:srgb2lin} gives
$Y_{\mathrm{ref}}=v^{2.4}$ and Eq.~\eqref{eq:slope} gives
$dY/dc=(2.4/1.055)v^{1.4}$, so
\begin{equation}
  \frac{Y_{\mathrm{ref}}}{dY/dc}
  = \frac{1.055\,v^{2.4}}{2.4\,v^{1.4}}
  = \frac{1.055\,v}{2.4}
  = \frac{c_{\mathrm{ref}}+0.055}{2.4},
  \label{eq:cancellation}
\end{equation}
and $2\times2.4=4.8$. Here $\mathrm{CSF}(u;Y_{\mathrm{ref}})$ denotes
Eq.~\eqref{eq:barten} evaluated with $E$ set from $Y_{\mathrm{ref}}$ by
Eqs.~\eqref{eq:pupil}--\eqref{eq:trolands}.

\paragraph{Equation~\eqref{eq:cancellation} is itself a result} and Sec.~10.3
depends on it: \textbf{the sRGB budget is linear in the code value, not in
the luminance.} The $2.4$ exponent cancels between the transfer function and
its slope, so gamma encoding has already performed most of the luminance
adaptation for free.

\paragraph{The measured size of the error.}
Over the full train split, legacy budget divided by calibrated budget:
\begin{center}
\begin{tabular}{@{}lccc@{}}
\toprule
 & low ($1.5\ \cpd$) & mid ($7.7\ \cpd$) & near-Nyquist ($37.5\ \cpd$) \\
\midrule
legacy / calibrated & $2.98\times$ & $1.97\times$ & $1.88\times$ \\
\bottomrule
\end{tabular}
\end{center}
decomposing at Nyquist as $3.35\times$ from $\mu=0.5$ alone and $0.56\times$
from $E=500$~Td pushing back, for a net $1.88\times$. \textbf{Every $\tau$
quoted under the legacy convention describes a patch roughly twice as
visible as claimed}, and roughly three times at low frequency. This is
finding \textbf{F2}.

% =====================================================================
\section{The residual: shape, scale, and what survives compositing}

\subsection{A smooth spectral reparameterisation, not a projection}

Given an unconstrained parameter $\mathbf{z}$, the residual is built so that
its spectrum respects $B$ \emph{by construction} \novel{}:
\begin{align}
  \mathbf{Z} &= \mathcal{F}\mathbf{z}, \label{eq:Z}\\
  \widehat{\mathbf{Z}}(f) &= B(f)\,
     \frac{\mathbf{Z}(f)}{\sqrt{1+|\mathbf{Z}(f)|^2}}, \label{eq:reparam}\\
  \boldsymbol{\delta}_0 &= \mathcal{F}^{-1}\widehat{\mathbf{Z}},
  \label{eq:delta0}
\end{align}
where $|\cdot|$ is complex modulus and the operations are elementwise over
bins $f\in\mathcal{K}$. Since $|\mathbf{Z}|/\sqrt{1+|\mathbf{Z}|^2}<1$ for
every finite $\mathbf{Z}$, we have $|\widehat{\mathbf{Z}}(f)|<B(f)$ for every
bin and every input --- an unconditional bound.

Three properties are load-bearing:
\begin{itemize}
\item \textbf{It is a reparameterisation, not a projection.} A hard
  projection zeroes the gradient of every component sitting on its bound,
  which is exactly where an adversarial component wants to live. The squash
  in Eq.~\eqref{eq:reparam} is smooth everywhere, the bound is never
  attained exactly, and no component can become stuck.
\item \textbf{Phase is preserved exactly.} The scaling multiplies the
  \emph{complex} coefficient, so only magnitude is limited. This also avoids
  differentiating $\arg(\mathbf{Z})$, which is undefined at the origin.
\item \textbf{It is scale-invariant in $\mathbf{z}$ up to shape.} Only the
  shape of the initialisation matters, which is why
  $\mathbf{z}^{(0)}\sim\mathcal{N}(\mathbf{0},\mathbf{I})$ rather than
  $\mathbf{0}$: Eq.~\eqref{eq:scale} below is $0/0$ at
  $\boldsymbol{\delta}_0=\mathbf{0}$.
\end{itemize}
Because $\mathcal{F}^{-1}$ of a half-spectrum is real by construction, no
conjugate-symmetry bookkeeping is required.

\subsection{The visibility index}
\label{sec:visibility}

Per-bin compliance is not enough, because the visual system sums evidence
across frequency and orientation channels: many individually sub-threshold
components can be jointly visible. This is \emph{probability summation},
standardly modelled by a Minkowski norm with exponent $\beta\approx3$--$4$
\citep{quick1974,robson1981,watson1979}. Define the per-bin visibility and
the pooled index
\begin{align}
  v(f) &= \kappa\,\big|\mathcal{F}\boldsymbol{\delta}(f)\big|\cdot
          \mathrm{CSF}\big(u(f)\big), \label{eq:vbin}\\
  V_\beta(\boldsymbol{\delta}) &=
     \left(\sum_{\mathrm{ch}=1}^{3}\ \sum_{f\in\mathcal{K}}
       v_{\mathrm{ch}}(f)^{\beta}\right)^{1/\beta},
     \qquad \beta = 3 .
  \label{eq:visibility}
\end{align}
Here $v(f)$ is the contrast of bin $f$ expressed in threshold units (so
$v=1$ is exactly at threshold), the inner sum runs over half-spectrum bins,
the outer sum over the three colour channels --- each channel is transformed
separately --- and $V_\beta\le1$ means at or below the detection threshold
of the modelled observer \emph{under the assumed geometry}.

\paragraph{Why not the max.}
$\max_f v(f)$ is the $\beta\to\infty$ limit and therefore the most
permissive criterion available; defaulting to it would flatter every result.
Concretely: a residual whose every bin respected its own budget measured a
max-visibility of $0.235$ while reaching a pixel-space deviation of
$\pm7.65$ --- nominally invisible and in fact catastrophic. Report
$V_\beta$; report $\max_f v(f)$ only as a secondary diagnostic.

\paragraph{The channel term is not cosmetic.}
Because the outer sum runs over three channels, a bound derived on one
channel is $3^{1/\beta}=1.44\times$ too generous. This mistake was made once
and caught by measurement (Sec.~6.2).

\subsection{Shape and scale are two separate operations}

The full reparameterisation composes the envelope with a global rescale
\novel{}:
\begin{equation}
  \boldsymbol{\delta}_1 = \boldsymbol{\delta}_0\cdot
     \frac{\tau}{V_\beta\!\left(\boldsymbol{\delta}_0\odot\mathcal{M}\right)} ,
  \label{eq:scale}
\end{equation}
where $\odot\mathcal{M}$ restricts to the silhouette actually pasted (for a
square patch $\mathcal{M}\equiv1$ and the restriction is vacuous), and the
scale is computed \emph{per sample} so that one image in a batch cannot
borrow another's headroom.

The division of labour is exact and worth stating in the thesis in these
terms:
\begin{center}
\begin{tabular}{@{}lL{9.6cm}@{}}
\toprule
\textbf{shape} & Eq.~\eqref{eq:reparam}: the $1/\mathrm{CSF}$ envelope biases
the residual's energy toward the bands the eye discards. \\
\textbf{scale} & Eq.~\eqref{eq:scale}: one global factor sets total perceived
visibility to exactly $\tau$. \\
\bottomrule
\end{tabular}
\end{center}
The two must be separate. The envelope alone bounds each bin but says
nothing about their sum; the rescale alone would let the attack put energy
wherever it likes, including where the eye is most sensitive. Normalising to
\emph{exactly} $\tau$ rather than clamping at it is deliberate: the attack
always wants the largest perturbation it is allowed, so the constraint is
active at every step anyway, and an equality is smooth where a $\min$ would
introduce a non-differentiable kink at the bound.

\paragraph{Corollary --- $\tau$ acts once, not twice.}
$\tau$ appears in both Eq.~\eqref{eq:budget} and Eq.~\eqref{eq:scale}, but
its occurrence in the budget is inert whenever $A_{\max}$ is inactive:
$B\propto\tau$ implies $\boldsymbol{\delta}_0\propto\tau$ implies
$V_\beta(\boldsymbol{\delta}_0)\propto\tau$, and the ratio in
Eq.~\eqref{eq:scale} cancels it. Measured at $S=128$ and the default
geometry, varying $\tau$ in Eq.~\eqref{eq:budget} alone over $[0.05,1]$
leaves the residual bit-identical.

\subsection{The range fit: why clipping cannot be used}
\label{sec:fitrange}

The residual is added to a base $\mathbf{r}$ and the sum must lie in
$[0,1]$. A clamp is the obvious choice and is fatal. A clamp is a hard
nonlinearity: it does not scale a residual, it \emph{bends} it, and a bent
peak is a corner, and a corner is broadband --- energy lands in exactly the
low and mid bands where $\mathrm{CSF}$ is largest and the budget tightest.
Measured: a residual built to exactly $1.0$ JND rose to $5.0$ JND after
clamping against a real image, with only $8\%$ of pixels actually clipped.

A \emph{uniform rescale} is the only operation that leaves the spectrum's
shape untouched --- scaling $\boldsymbol{\delta}$ by $c$ scales every
Fourier coefficient by $c$, so a CSF-shaped residual stays CSF-shaped and
$V_\beta(c\boldsymbol{\delta})=c\,V_\beta(\boldsymbol{\delta})$ exactly.
Define per-pixel headroom and the admissible ratio \novel{}
\begin{align}
  h_i &= \begin{cases}
      1-r_i, & \delta_i>0 \quad\text{(pushing up)},\\
      r_i,   & \delta_i\le0 \quad\text{(pushing down)},
  \end{cases} \label{eq:headroom}\\[3pt]
  \gamma_i &= \begin{cases}
      +\infty, & h_i \le h_{\min}\ \text{ or } i\notin\mathrm{supp}(\mathcal{M}),\\[2pt]
      \dfrac{h_i}{|\delta_i|}, & \text{otherwise},
  \end{cases} \label{eq:ratio}\\[3pt]
  \boldsymbol{\delta}_2
     &= \boldsymbol{\delta}_1 \cdot
        \mathrm{clip}\!\left(Q_q\big(\{\gamma_i\}\big),0,1\right),
  \label{eq:fit}
\end{align}
where $i$ indexes pixels and channels, $r_i$ is the base value at $i$,
$Q_q$ is the empirical $q$-quantile with $q=0.001$, and
$h_{\min}=10^{-3}$.

Two exclusions in Eq.~\eqref{eq:ratio} are each worth a paragraph in the
thesis because each caused a silent total failure:
\begin{itemize}
\item \textbf{Saturated pixels are excluded, not down-weighted.} A pixel
  already at $0$ or $1$ has zero headroom in one direction, so its ratio is
  $0$ and it drags any quantile to $0$ with it. Over a $128\times128$ patch
  a $0.1\%$ quantile tolerates $\approx49$ such pixels; over a full
  $512\times1024$ frame it tolerates $524$, and a single Cityscapes sky
  blow-out exceeds that. The observed consequence was a full-image probe
  that scaled \emph{every} band to exactly zero and produced no measurement.
  Excluding them is also physically correct: where the base is already $1.0$
  and the residual pushes up, the pixel was never going to move, so it
  should not veto the scale of every pixel that could.
\item \textbf{The mask is not optional.} A shaped patch loads its cutout
  padded with transparent black, so for one reference $45\%$ of the square
  is RGB $0$ and $95\%$ of the region outside the silhouette is. Unmasked,
  those never-pasted pixels drove the quantile to $0.0000$ and scaled the
  residual to exactly zero --- a five-epoch run that produced the unmodified
  reference image. The same quantile computed inside the silhouette was
  $1.26$, i.e.\ the residual would have fitted at full scale with room to
  spare.
\end{itemize}
Allowing the lowest $q$ of pixels to clip trades a tiny, measurable
violation for a usable attack. The violation is not assumed away: it is
measured by Sec.~5.5.

\subsection{Enforcement: bounding what survives, not what was intended}
\label{sec:enforce}

Everything above constrains the residual we \emph{intend} to add. The
observer sees
\begin{equation}
  \boldsymbol{\delta}^{\mathrm{real}}
  = \mathrm{clip}(\mathbf{r}+\boldsymbol{\delta}_2,0,1) - \mathbf{r},
  \label{eq:realresid}
\end{equation}
which is a different signal, and the difference is not small. The honest
number to report is therefore the \emph{realised visibility}
\begin{equation}
  V^{\mathrm{real}}
  = V_\beta\!\left(\boldsymbol{\delta}^{\mathrm{real}}\odot\mathcal{M}\right).
  \label{eq:realvis}
\end{equation}
$\tau$ is an intent; $V^{\mathrm{real}}$ is an outcome; the gap between them
is the measure of how well the constraint held.

\paragraph{The gap was large (finding F3).}
On SegFormer, image 420, $\tau=0.25$, five seeds at 1000 steps:
$V^{\mathrm{real}}$ came out $2.77\times\tau$ on the $\mu=0.5$ convention
(worst seed $4.52\times$) and $7.44\times$ on the locally measured
convention (worst $12.34\times$). The constraint held at initialisation
($\tau=0.25$ requested, $0.2508$ realised) and broke \emph{during}
optimisation, because the attack pushes toward the largest perturbation it
is allowed and so drives more pixels into the clamp. The signature is
diagnostic: the worst seed for visibility had the \emph{lowest} residual RMS
and absolute maximum of the five --- backwards for ``bigger is more
visible'', and exactly right for clipping, which removes peaks while adding
harmonics.

\paragraph{The fix is the one operation that cannot distort a spectrum.}
Bisect on a single scalar $c\in[0,1]$ and return $c\,\boldsymbol{\delta}_2$,
choosing the largest $c$ whose \emph{composited} residual measures within
$\tau$:
\begin{equation}
  c^\star = \max\Big\{c\in\mathcal{C} :
    V_\beta\Big(\big[\mathrm{clip}(\mathbf{r}+c\,\boldsymbol{\delta}_2,0,1)
    -\mathbf{r}\big]\odot\mathcal{M}\Big)\le\tau\Big\},
  \label{eq:bisect}
\end{equation}
where $\mathcal{C}$ is the finite set of candidates the bisection actually
evaluates ($12$ iterations by default).

\paragraph{The soundness argument, which matters because $V$ need not be
monotone in $c$.} The bisection maintains an invariant: the lower bracket is
only ever assigned a value that has been \emph{measured} compliant, the
upper bracket only ever one measured violating, and the lower bracket is
what is returned. The bound therefore holds regardless of monotonicity ---
and monotonicity is \emph{not} guaranteed, precisely because clipping
generates harmonics. A bisection assuming it would be unsound here; this one
is not, because it never returns a value it has not checked. $c=0$ is
trivially compliant since $V_\beta(\mathbf{0})=0$, so the lower bracket
always starts valid and the search cannot fail.

\paragraph{Gradient contract.}
$c^\star$ is computed under \texttt{no\_grad} and treated as a constant, so
the backward pass is the gradient of the scaled residual --- the same
contract as gradient clipping. The forward pass is exactly compliant.

\paragraph{Cost.}
$12$ transforms on the patch grid per step and no model forward, which is
negligible beside one SegFormer pass. The attack does get weaker, and that
is the point: the overshoot was buying effectiveness with perceptual budget
the run had already promised not to spend. Measured cost: $4.58$ mIoU
points (Sec.~10.4).

\subsection{The complete per-image pipeline}

Composing Secs.~5.1--5.5, the patch presented to the network is
\begin{equation}
  \mathbf{P} = \mathrm{clip}\Big(\mathbf{r} +
    \underbrace{c^\star\cdot
      \underbrace{Q_q(\gamma)\cdot
        \underbrace{\tfrac{\tau}{V_\beta(\boldsymbol{\delta}_0\odot\mathcal{M})}
          \underbrace{\mathcal{F}^{-1}\!\left[
            \tfrac{B\,\mathcal{F}\mathbf{z}}
                  {\sqrt{1+|\mathcal{F}\mathbf{z}|^2}}\right]}
            _{\text{shape}}}_{\text{scale to }\tau}}
      _{\text{range fit}}}_{\text{enforce realised }\tau},\ 0,\ 1\Big).
  \label{eq:pipeline}
\end{equation}
The order is load-bearing: enforcement comes \emph{after} the range fit, not
instead of it. The range fit keeps the composite inside $[0,1]$ for all but a
quantile of pixels; enforcement bounds what the remaining clipping
\emph{costs}. Both are uniform rescales, so the spectrum's shape survives
either way.

\paragraph{The base is the covered content.} \novel{}
With \texttt{--from\_image}, $\mathbf{r}$ is the image region the patch will
occupy, so the patch is that content plus a bounded residual rather than an
arbitrary image. This is what makes ``the patch looks like what was already
there'' true by construction rather than by a penalty term, and it is the
structural difference from \citet{tan2021lap}, whose $L_{\mathrm{rat}}$ pulls
toward a reference \emph{softly} and can therefore be outvoted by the
adversarial term.

% =====================================================================
\section{The universal variant: one residual for the whole dataset}

\subsection{What changes, and what deliberately does not}

The per-image mode of Sec.~5 optimises a fresh residual per image and takes
its base from that image. The universal mode differs in \emph{exactly two}
ways and in nothing else \citep[cf.][]{moosavi2017universal}:
\begin{enumerate}
\item $\boldsymbol{\delta}$ is \emph{shared} across the whole dataset;
\item the CSF is evaluated at one \emph{fixed} reference luminance, because
  a universal patch cannot look at the image it lands on.
\end{enumerate}
It is the non-adaptive control for the per-image attack, and its value is as
much in a legible null result as in a success.

The composite is
\begin{equation}
  \mathbf{x}^{\mathrm{adv}}_{\mathrm{win}}
  = \mathrm{clip}\big(\mathbf{x}^{01}_{\mathrm{win}}
    + \boldsymbol{\delta},\,0,\,1\big),
  \label{eq:univcomposite}
\end{equation}
where $\mathbf{x}^{01}_{\mathrm{win}}$ is the footprint window of the image
in $[0,1]$ code values. Every image receives the \emph{same residual on
different content}, which is the whole point.

\subsection{Hard projection, and why the envelope must be normalised}
\label{sec:univproj}

Here the constraint is a true projection rather than a squash, because the
bound must hold as an \emph{inequality on the parameter itself} for
``project the parameter, not the render'' to be checkable:
\begin{align}
  \widehat{\mathbf{Z}}(f)
    &= \mathbf{Z}(f)\cdot\min\!\left(1,\
       \frac{B'(f)}{|\mathbf{Z}(f)|}\right), \label{eq:project}\\
  \pi &= \frac{1}{|\mathcal{K}|}\Big|\Big\{f:\
       |\mathbf{Z}(f)|>B'(f)>0\Big\}\Big| ,
  \label{eq:atbound}
\end{align}
with $\pi$ the fraction of live bins sitting at their bound. Phase is again
untouched. The dead gradient at a saturated bin is real, so $\pi$ is
reported rather than discarded: a run whose $\pi$ climbs toward $1$ has
stopped optimising and is only re-phasing.

\paragraph{Why $B'$ and not $B$.}
The obvious construction --- project onto the raw $1/\mathrm{CSF}$ budget,
then rescale the result so that pooled visibility equals $\tau$ --- does not
work. The projection clamps most bins down, so pooled visibility afterwards
is far below $\tau$, and the rescale that restores it multiplies every bin
back up: a smoke test put bins at $90\times$ their own budget. Per-bin bound
and pooled equality are not simultaneously satisfiable on one envelope.

Normalise the \emph{envelope} instead \novel{}:
\begin{align}
  V_{\mathrm{full}} &= \left(3\sum_{f\in\mathcal{K}}
     \Big(\kappa\,B(f)\,\mathrm{CSF}(u(f))\Big)^{\beta}\right)^{1/\beta},
     \label{eq:vfull}\\
  B'(f) &= B(f)\cdot\frac{\tau}{V_{\mathrm{full}}},
     \label{eq:normbudget}
\end{align}
where $V_{\mathrm{full}}$ is the pooled visibility of a hypothetical
residual sitting at budget on \emph{every} bin, and the factor $3$ is the
colour-channel count of Eq.~\eqref{eq:visibility}. Then the per-bin bound is
the real constraint and pooled visibility is an \emph{inequality},
$V_\beta\le\tau$, with equality only if every bin saturates --- the honest
direction for the inequality to point.

\paragraph{The factor 3 was worth measuring.}
Omitting it makes the envelope $3^{1/3}=1.44\times$ too generous. Measured
before the term was added: pooled $0.0541$ against a nominal $\tau=0.05$, an
$8\%$ overshoot that every assertion in the training loop passed, because
the assertion and the normaliser shared the mistake.

\paragraph{Consequence for comparison.}
A universal run at nominal $\tau$ is \emph{less} visible than a per-image
run at the same nominal $\tau$, where Eq.~\eqref{eq:scale} enforces equality.
The two must therefore be compared at matched \emph{realised} visibility as
well as at matched nominal $\tau$.

\subsection{Projection belongs in the optimiser step, not the forward pass}

Placing Eq.~\eqref{eq:project} inside the forward pass makes the map
radially flat above the bound, so $\partial\mathcal{L}/\partial|\mathbf{Z}|$
is \emph{exactly} zero at any saturated bin and its magnitude can never come
back down. Saturation becomes an absorbing state. Measured from the default
initialisation: $\pi$ sat at $0.989$ for 200 Adam steps while the per-bin
spend ratio moved by $2\times10^{-5}$ --- the optimiser could only rotate
phase, and the spectral allocation this mode exists to measure was frozen at
$1.000$ by construction rather than chosen. Applying the projection as a
post-step parameter clamp, with the loss seeing an unflattened gradient,
fixes it.

\subsection{Geometry is pinned}

The universal mode refuses to run unless $S=p$. Resampling a residual
resamples its spectrum, so the budget its bins were projected onto would
stop describing the pasted signal, and $\tau$ would silently cease to mean
anything. At $H=512$, $s=0.25$ that is $S=p=128$. The same warning applies
to the per-image mode whenever \texttt{--patch\_scale} and
\texttt{--patch\_size} are changed independently.

\subsection{Calibrated visibility: the headline metric for this mode}

$\tau$ is set once at a fixed $Y_{\mathrm{ref}}$. What that one residual
actually \emph{costs} on each image is
\begin{align}
  Y_i &= \frac{\sum_{(u,v)\in\mathcal{M}} Y\big(\mathbf{x}^{01}_{i,uv}\big)}
              {|\mathcal{M}|}, \label{eq:Yi}\\
  \tau^{\mathrm{cal}}_i &= V_\beta\Big(\boldsymbol{\delta}\odot\mathcal{M};\
     \mathrm{CSF}(\cdot\,;Y_i),\ \kappa_{\mathrm{cal}}(Y_i)\Big),
  \label{eq:taucal}
\end{align}
with $Y(\cdot)$ from Eq.~\eqref{eq:luma}, $\mathrm{CSF}(\cdot;Y_i)$ from
Eq.~\eqref{eq:barten} with $E$ set by
Eqs.~\eqref{eq:pupil}--\eqref{eq:trolands}, and $\kappa_{\mathrm{cal}}$ from
Eq.~\eqref{eq:kappacal}. \emph{Both} the sensitivity and the contrast
denominator respond to $Y_i$, so this is not a rescale of the nominal number
and the per-image ordering can differ from it.

The sentence this exists to license is: \emph{``the budget was set at
$\tau=a$ using a fixed reference; realised visibility across the val set
ranged from $b$ to $c$.''} That spread is the empirical case for content
dependence, or the evidence against it.

\subsection{Clipping is tracked, not assumed away}

Equation~\eqref{eq:univcomposite} truncates where content is near $0$ or
$1$, so the realised perturbation is not the projected one and $\tau$ is
violated in the permissive direction. Report
\begin{equation}
  \mathrm{frac_{clip}}
  = \frac{1}{3|\mathcal{M}|}
    \Big|\Big\{i\in\mathcal{M}\times\{1,2,3\}:\
    x_i+\delta_i\notin[0,1]\Big\}\Big|
  \label{eq:fracclip}
\end{equation}
every step. A non-negligible value invalidates the guarantee quietly, which
is the worst way for it to fail. The alternative composite
(\texttt{--csf\_composite fit}, Eq.~\eqref{eq:fit}) rescales rather than
truncates and so preserves the spectrum --- at the cost of a per-image
scale, which is precisely the per-image adaptation a universal patch is not
supposed to have. Clipping is the default for that reason. Note also that
\texttt{clip} has zero gradient on the saturated side, so a heavily clipped
run also loses gradient on those pixels.

% =====================================================================
\section{Application, threat model, and the optimisation loop}

\subsection{Compositing}

Following the localised-patch threat model of \citet{brown2017patch}, with
optional silhouette mask after \citet{tan2021lap},
\begin{align}
  \tilde{\mathbf{P}} &= \big(\mathrm{Resize}_{p\times p}(\mathbf{P})
     - \boldsymbol{\mu}_n\big)\oslash\boldsymbol{\sigma}_n,
     \label{eq:normalise}\\
  \mathbf{m} &= \mathrm{Pad}_{t,\ell}\big(\mathbb{1}_{p\times p}
     \odot\mathcal{M}\big), \label{eq:footprint}\\
  \mathbf{x}^{\mathrm{adv}} &= \mathbf{x}\odot(1-\mathbf{m})
     + \mathrm{Pad}_{t,\ell}(\tilde{\mathbf{P}})\odot\mathbf{m},
     \label{eq:apply}
\end{align}
where $\mathbf{P}$ is the rendered patch in $[0,1]$ from
Eq.~\eqref{eq:pipeline}, $\oslash$ is elementwise division, $(t,\ell)$ is
the top-left placement, and $\mathrm{Pad}_{t,\ell}$ zero-pads the patch into
the full frame at that offset. The footprint follows the \emph{silhouette},
not the bounding box, so with a cutout only genuinely occluded pixels are
excluded from scoring --- stricter and more honest than excluding the whole
square.

\paragraph{Placement.}
Three modes: \texttt{center}; \texttt{semantic}, which reads the clean
prediction; and \texttt{gradcam}, which reads a sensitivity map. The last
two impose an ordering constraint --- clean forward, then resolve placement,
then apply --- and getting it backwards silently falls back to centre, so
that every downstream number is quietly about a different patch position.

\subsection{Optimisation}

\begin{equation}
  \mathbf{z}^{(0)}\sim\mathcal{N}(\mathbf{0},\mathbf{I}),\qquad
  \mathbf{z}^{(j+1)} = \mathrm{Adam}\big(\mathbf{z}^{(j)},
     \nabla_{\mathbf{z}}\mathcal{L}\big),\qquad j=0,\dots,K-1,
  \label{eq:opt}
\end{equation}
with $\mathcal{L}$ one of the objectives of Sec.~8, $K$ the step count, and
Adam \citep{kingma2015adam} on the raw gradient. The network $f$ is frozen
throughout and $\nabla_\phi f$ is never formed. The learning rate follows a
cosine schedule \citep{loshchilov2017sgdr}.

\paragraph{The learning rate is not transferable between modes, and this is
finding F6.} The two modes optimise different objects:
\begin{center}
\small
\begin{tabular}{@{}lL{5.1cm}L{4.9cm}@{}}
\toprule
 & \textbf{per-image (\texttt{csf})} & \textbf{universal (\texttt{universal\_csf})} \\
\midrule
what $\mathbf{z}$ is & a \emph{direction}: mapped through a scale-invariant
squash and renormalised to exactly $\tau$ & \emph{the residual itself}, in
pixel units (RMS $0.027$ at $\tau=0.25$) \\
what lr controls & how fast that direction rotates & movement inside a small
ball; overshoot is discarded by the projection \\
correct lr & $0.2$ & $0.01$ \\
\bottomrule
\end{tabular}
\end{center}
Scaling the per-image parameter by $1000\times$ moves the rendered patch by
$7.5\%$. Measured mean relative change of the applied residual per step:
\begin{center}
\begin{tabular}{@{}lcc@{}}
\toprule
lr & \texttt{csf} & \texttt{universal\_csf} \\
\midrule
$0.2$   & $0.50$   & $1.89$ \\
$0.05$  & $0.084$  & $1.09$ \\
$0.01$  & $0.031$  & $0.15$ \\
$0.001$ & $0.0038$ & $0.039$ \\
\bottomrule
\end{tabular}
\end{center}
A relative change above $1$ means each step discards the previous residual
entirely. Corroborating diagnostic: $\pi$ after 200 steps is $0.931$ at
lr~$0.2$ (pinned to the boundary, phase-only) against $0.508$ at lr~$0.01$
(healthy interior). \textbf{Watch $\pi$ in the logs}: pinned above
$\approx0.9$ means the learning rate is too large whatever any script says.

\paragraph{Annealing is not optional at these rates.}
A flat learning rate was measured swinging $9.6$ mIoU points across four
\emph{identical} single-image invocations, because Adam's step is
$\approx\mathrm{lr}$ per coordinate whatever the gradient is, so the run
never settles and the reported number is wherever the walk happened to stop.

% =====================================================================
\section{The three untargeted objectives}
\label{sec:losses}

All three objectives below are \emph{untargeted}: they specify no class the
attack should impose, only that the prediction should be wrong or uncertain.
They are the axis on which the loss comparison is run. The targeted
objective (IPatch-style, driving every scored pixel toward one class
\citealp{mirsky2023ipatch}) is a separate arm and is not part of this axis;
it is mentioned only where the contrast is informative.

\subsection{The scored support, shared by all three}
\label{sec:support}

Every objective is a mean over a pixel set $\Omega$, never over the whole
frame:
\begin{equation}
  \Omega = \underbrace{\{(u,v):\mathbf{y}(u,v)\neq255\}}_{\text{valid labels}}
  \ \setminus\ \underbrace{\mathrm{supp}(\mathbf{m})}_{\text{patch footprint}}
  \ \cap\ \underbrace{\mathcal{R}}_{\text{reach mask (optional)}} .
  \label{eq:support}
\end{equation}

\paragraph{Footprint exclusion is not optional.}
Leave the patch region in the loss and the cheapest way to reduce it is to
make the \emph{patch} look wrong, which collapses the objective while the
remote effect stays exactly zero \citep{mirsky2023ipatch}. The attack we are
studying is a \emph{remote} one; scoring the footprint measures occlusion.

\paragraph{Reach restriction, and what it can and cannot buy.}
A patch can only influence pixels inside the model's effective receptive
field \citep{luo2016erf}. Measured with a label-free random-stimulus probe,
prediction-change collapses below $1\%$ past $\approx300$~px at
$512\times1024$ and $\approx150$~px at $256\times512$; of $\approx450$k
remote pixels only $\approx50$--$80$k are reachable. The unreachable
$\approx85\%$ contribute gradients that are small but not zero, and that do
not correspond to achievable flips: they dilute the mean by $\approx5\times$
and add variance to the update direction. Restricting $\Omega$ to
$\mathcal{R}$ therefore (a) shrinks the denominator and (b) removes those
contributions outright. \textbf{It cannot extend reach} --- the reachable
set is fixed by the architecture. Available as a radial disc of default
radius $0.62H$ (set from the measured collapse point, not chosen) or as the
measured mask itself.

\paragraph{Reporting asymmetry, to be declared.}
The loss support and the evaluation metric cover different pixel sets: we
optimise on $\mathcal{R}$ but still report the drop over \emph{all} remote
pixels. This is deliberate --- the metric must not be tuned to the method
--- and must be stated wherever a reach-restricted number appears.

\paragraph{The guarded mean.}
Every objective reduces with
\begin{equation}
  \langle g\rangle_\Omega
  = \frac{\sum_{(u,v)} g_{uv}\,\mathbb{1}[(u,v)\in\Omega]}
         {\max\big(|\Omega|,\,1\big)} ,
  \label{eq:reduce}
\end{equation}
where $g_{uv}$ is the per-pixel quantity. The clamped denominator is not
defensive programming for its own sake: Cityscapes batches occur in which
void regions (ego-vehicle hood, borders) plus the patch footprint cover
every pixel, giving $0/0$. The resulting NaN propagates through the backward
pass, permanently poisons Adam's moment estimates, and freezes the patch at
its initial value for the rest of training with no error raised. That cost
one 100-epoch run before it was noticed.

\subsection{Objective 1 --- cross-entropy}

The simplest untargeted objective maximises the network's cross-entropy
against the ground truth:
\begin{align}
  \mathrm{CE}_{uv} &= -\log p_{uv,\,\mathbf{y}(u,v)},
     \label{eq:ce}\\
  \mathcal{L}_{\mathrm{CE}}
     &= -\big\langle \mathrm{CE}_{uv}\big\rangle_\Omega
      = \big\langle \log p_{uv,\,\mathbf{y}(u,v)}\big\rangle_\Omega .
  \label{eq:celoss}
\end{align}
Here $p_{uv,k}$ is the softmax probability of class $k$ at pixel $(u,v)$
from Sec.~2, and $\mathbf{y}(u,v)$ is that pixel's ground-truth class.
Cross-entropy is \emph{high} when the prediction is wrong, so the attack
maximises it; since Adam minimises, the returned scalar carries a leading
minus sign. This is the segmentation analogue of the untargeted FGSM/PGD
objective \citep{goodfellow2015fgsm,madry2018pgd}.

\paragraph{Its known pathology, which motivates the other two.}
$-\log p_y$ is unbounded above, so once a pixel is thoroughly misclassified
it can still contribute a large and growing gradient. The mean is then
dominated by a few already-won pixels, and the attack keeps spending
capacity where it has already succeeded instead of on the pixels still
holding out. \citet{gu2022segpgd} make the same observation for segmentation
PGD and address it by reweighting correctly- and incorrectly-classified
pixels over the attack schedule.

\subsection{Objective 2 --- CosPGD}

CosPGD \citep{agnihotri2024cospgd} rescales the per-pixel cross-entropy by
the cosine similarity between the predicted distribution and the one-hot
ground truth:
\begin{align}
  w_{uv} &= \cos\!\big(\mathbf{p}_{uv},\,\mathbf{e}_{\mathbf{y}(u,v)}\big)
     = \frac{\langle\mathbf{p}_{uv},\mathbf{e}_{\mathbf{y}(u,v)}\rangle}
            {\|\mathbf{p}_{uv}\|_2\,\|\mathbf{e}_{\mathbf{y}(u,v)}\|_2}
     = \frac{p_{uv,\,\mathbf{y}(u,v)}}{\|\mathbf{p}_{uv}\|_2},
     \label{eq:cosw}\\
  \mathcal{L}_{\mathrm{CosPGD}}
     &= -\big\langle \sg{w_{uv}}\cdot\mathrm{CE}_{uv}\big\rangle_\Omega ,
  \label{eq:cospgd}
\end{align}
where $\mathbf{e}_k$ is the one-hot vector of class $k$ (so
$\|\mathbf{e}_k\|_2=1$, giving the third equality), and
$\sg{\cdot}$ denotes a stop-gradient.

\paragraph{The cosine is a weight, never the objective.}
$w_{uv}\in(0,1]$ measures how aligned the prediction still is with the
truth. Pixels still predicted correctly ($w\approx1$) keep full weight;
pixels already fooled ($w\approx0$) are de-emphasised. The gradient itself
comes from CE, which does not saturate the way a raw cosine does near
alignment. This is the direct answer to Sec.~8.2's pathology: effort is
concentrated on the pixels that have not yet flipped.

\paragraph{Two declared deviations from the paper.}
\begin{enumerate}
\item \textbf{The cosine is detached.} \citet{agnihotri2024cospgd} do not
  detach it; in their Algorithm~1 the gradient flows through the whole
  product. Detaching removes gradient contributions through the weight and
  is a deliberate deviation, common in patch-attack implementations.
\item \textbf{Adam replaces sign-SGD with $\varepsilon$-projection.} CosPGD
  is a PGD variant \citep{madry2018pgd}, and a patch is not an
  $\varepsilon$-ball perturbation, so Eq.~\eqref{eq:cospgd} is optimised
  with Adam on the raw gradient.
\end{enumerate}
Both are correct for this threat model and both must be stated in the
thesis, because a reader comparing numbers against the CosPGD paper is
otherwise comparing two different algorithms.

\subsection{Objective 3 --- Tsallis entropy}
\label{sec:tsallis}

\textbf{Status: specified here, not yet implemented.} This subsection is the
implementation contract as well as the derivation. The entropy itself is
Tsallis's \citeyearpar{tsallis1988} one-parameter generalisation of Shannon
entropy \citep{shannon1948}, independently introduced by
\citet{havrda1967}; its use as an \emph{adversarial objective for
segmentation} carries no external source here and is marked \novel{}.

\paragraph{Definition.}
For a distribution $\mathbf{p}\in\Delta^{C-1}$ and an entropic index
$q>0$, $q\neq1$,
\begin{equation}
  S_q(\mathbf{p}) = \frac{1}{q-1}\left(1-\sum_{c=1}^{C}p_c^{\,q}\right).
  \label{eq:tsallis}
\end{equation}
$q$ is the single hyper-parameter; $C$ is the class count; $p_c$ is the
probability of class $c$. $S_q\ge0$ with equality if and only if
$\mathbf{p}$ is one-hot, and $S_q$ is maximised by the uniform distribution
$p_c=1/C$, where it attains the \emph{$q$-logarithm} of $C$:
\begin{equation}
  S_q^{\max} = \ln_q C = \frac{C^{1-q}-1}{1-q} .
  \label{eq:lnq}
\end{equation}
For $C=19$ this is $6.718$ at $q=0.5$, $\ln 19=2.944$ in the Shannon limit,
and $0.947$ at $q=2$ --- the range varies by an order of magnitude with $q$,
which is why the objective must be normalised before different $q$ are
compared:
\begin{equation}
  \bar S_q(\mathbf{p}) = \frac{S_q(\mathbf{p})}{\ln_q C}\ \in[0,1].
  \label{eq:tsallisnorm}
\end{equation}

\paragraph{The objective.}
\begin{equation}
  \mathcal{L}_{\mathrm{Tsallis}}
  = -\big\langle \bar S_q(\mathbf{p}_{uv})\big\rangle_\Omega ,
  \label{eq:tsallisloss}
\end{equation}
with $\Omega$ from Eq.~\eqref{eq:support} and $\langle\cdot\rangle_\Omega$
from Eq.~\eqref{eq:reduce}. Minimising $\mathcal{L}_{\mathrm{Tsallis}}$
maximises the mean normalised entropy of the remote prediction: the attack
drives the network toward \emph{indecision} rather than toward a specific
error. It is the sign-flipped counterpart of entropy \emph{minimisation} as
used in semi-supervised learning and test-time adaptation
\citep{grandvalet2004,wang2021tent}.

\paragraph{The Shannon limit.}
Writing $q=1+\epsilon$ and expanding $p_c^{\,1+\epsilon}
= p_c e^{\epsilon\ln p_c} = p_c(1+\epsilon\ln p_c) + O(\epsilon^2)$,
\begin{equation}
  S_{1+\epsilon}(\mathbf{p})
  = \frac{1}{\epsilon}\Big(1-\sum_c p_c-\epsilon\sum_c p_c\ln p_c\Big)
    + O(\epsilon)
  \ \xrightarrow[\epsilon\to0]{}\ -\sum_c p_c\ln p_c
  = H(\mathbf{p}),
  \label{eq:shannonlimit}
\end{equation}
using $\sum_c p_c=1$. So $q\to1$ recovers Shannon entropy exactly, and
Eq.~\eqref{eq:tsallisnorm} recovers the normalised entropy
$H(\mathbf{p})/\ln C$ that the diagnostics already report. \textbf{An
implementation must dispatch $q=1$ to the Shannon form}, since
Eq.~\eqref{eq:tsallis} divides by $q-1$.

\paragraph{Gradient, in closed form.}
Let $\mathbf{z}_{uv}\in\mathbb{R}^C$ be the logits at a pixel and
$\mathbf{p}=\mathrm{softmax}(\mathbf{z})$. From Eq.~\eqref{eq:tsallis},
$\partial S_q/\partial p_c = -q\,p_c^{\,q-1}/(q-1)$, and the softmax
Jacobian is $\partial p_c/\partial z_k = p_c(\delta_{ck}-p_k)$ with
$\delta_{ck}$ the Kronecker delta. Chaining,
\begin{align}
  \frac{\partial S_q}{\partial z_k}
  &= \sum_{c}\frac{-q\,p_c^{\,q-1}}{q-1}\;p_c(\delta_{ck}-p_k)
   = \frac{-q}{q-1}\Big(p_k^{\,q} - p_k\sum_c p_c^{\,q}\Big)
  \nonumber\\
  &= \frac{q}{q-1}\;p_k\left(\sum_{c}p_c^{\,q} - p_k^{\,q-1}\right).
  \label{eq:tsallisgrad}
\end{align}
Reading it: a logit is pushed \emph{up} when its class is less probable than
the ``typical'' class under the $q$-weighting $\sum_c p_c^{\,q}$, and
\emph{down} when it is more probable. The objective is therefore a
\emph{levelling} force, and it needs no label to compute --- which is the
whole point.

\paragraph{Sanity check against the Shannon limit.}
Substituting $q=1+\epsilon$ into Eq.~\eqref{eq:tsallisgrad} and using
$\sum_c p_c^{1+\epsilon}\approx1+\epsilon\sum_c p_c\ln p_c$ and
$p_k^{\,\epsilon}\approx1+\epsilon\ln p_k$ gives
$\partial H/\partial z_k = -p_k\big(\ln p_k + H(\mathbf{p})\big)$, the
standard Shannon-entropy gradient. The closed form is therefore consistent
at the boundary.

\paragraph{What $q$ controls, and why it is worth sweeping.}
The factor $p_k^{\,q-1}$ in Eq.~\eqref{eq:tsallisgrad} sets which part of
the distribution the objective attends to.
\begin{center}
\small
\begin{tabular}{@{}lL{10.2cm}@{}}
\toprule
$q>1$ & Emphasises the \emph{dominant} class: $p_c^{\,q}$ suppresses the
tail, so the force is concentrated on pulling the winner down. Acts like an
attack on confidence. As $q\to\infty$, $S_q$ is governed by
$\max_c p_c$ alone. \\[3pt]
$q\to1$ & Shannon entropy: every class contributes in proportion to
$p_c\ln p_c$. \\[3pt]
$q<1$ & Emphasises the \emph{tail}: $p_c^{\,q}$ inflates small
probabilities, so the force pushes probability mass out into many classes.
As $q\to0^{+}$, $S_q\to C-1$ for any full-support $\mathbf{p}$ and the
objective degenerates into counting the support. \\
\bottomrule
\end{tabular}
\end{center}
$q$ is therefore a second realism-independent knob and a genuine axis:
$q\in\{0.5,\,1,\,2\}$ is the minimal informative sweep, run at fixed $\tau$.

\paragraph{Properties that make it worth having as a third objective.}
\begin{itemize}
\item \textbf{Label-free.} Equation~\eqref{eq:tsallisloss} never references
  $\mathbf{y}$. It can therefore be scored on void pixels, it cannot be
  accused of using ground truth a real attacker would not have, and it is
  the natural objective for the universal patch of Sec.~6, where no
  per-image label is available at attack time.
\item \textbf{Bounded.} $\bar S_q\in[0,1]$, so unlike Eq.~\eqref{eq:celoss}
  no pixel can dominate the mean without bound. This is the same benefit
  CosPGD obtains by reweighting, but obtained by construction.
\item \textbf{It targets a different failure mode.} CE and CosPGD produce
  \emph{confident errors}; a maximum-entropy attack produces
  \emph{abstention}. For a downstream consumer that thresholds on
  confidence, these are not the same attack at all, and the distinction is
  worth a paragraph in the discussion.
\end{itemize}

\paragraph{Two cautions that must be stated with any Tsallis result.}
\begin{enumerate}
\item \textbf{Maximum entropy is not maximum damage.} The uniform
  distribution is the global optimum of Eq.~\eqref{eq:tsallisloss}, and
  while reaching it implies an $\arg\max$ change in general, entropy can
  also rise substantially \emph{without} any $\arg\max$ changing. Never
  quote an entropy objective's success by its own loss value; quote
  $\Delta\mathrm{mIoU}_{\mathrm{rem}}$ and $\phi$ (Sec.~9), exactly as for
  the other two.
\item \textbf{Objective/metric coupling.} The diagnostic suite already
  reports normalised Shannon entropy by distance ring. That diagnostic is
  \emph{the objective} for $q\to1$, so it stops being independent evidence
  for a Tsallis run and becomes a training curve. Report it, label it as
  such, and carry the mechanism claim on the confusion flows and the reach
  curve instead.
\end{enumerate}

\paragraph{Numerical requirements for the implementation.}
$p_c^{\,q-1}$ diverges as $p_c\to0$ for $q<1$, so clamp
$p_c\ge\varepsilon=10^{-12}$ before exponentiating; compute $p_c^{\,q}$ as
$\exp(q\log p_c)$ on the clamped values for stability; dispatch $q=1$ to the
Shannon branch; and normalise by $\ln_q C$ with the \emph{same} $C$ used to
build $\Omega$. Reuse Eq.~\eqref{eq:reduce} verbatim rather than
\texttt{reduction="mean"}, for the reason given in Sec.~8.1.

\subsection{What the three objectives are for}

\begin{center}
\small
\begin{tabular}{@{}lL{3.0cm}L{3.2cm}L{3.9cm}@{}}
\toprule
 & \textbf{CE} & \textbf{CosPGD} & \textbf{Tsallis} \\
\midrule
uses labels & yes & yes & \textbf{no} \\
per-pixel range & unbounded & unbounded, reweighted & bounded in $[0,1]$ \\
effort allocation & drifts to won pixels & held on unflipped pixels & held on confident pixels \\
failure mode induced & confident error & confident error & abstention \\
role in the thesis & baseline & main objective & label-free control and a different threat \\
\bottomrule
\end{tabular}
\end{center}

\paragraph{The finding that applies to all three (F7).}
An untargeted objective has no class goal, yet it converges on \emph{one}
confusion channel. Measured: $\approx97\%$ of all remote flips went
road~$\to$~car, with car IoU dropping $\approx68$ points in every image. The
untargeted loss is an \emph{implicit class selector}, not a broad
disruptor --- it finds the most exploitable boundary and exploits only that.
Aggregate mIoU hides this completely; the source$\to$destination flow table
is what makes it visible, and a single flow carrying more than half of all
flips is the criterion for calling a result \emph{structured} rather than
\emph{diffuse}. These have different security implications and must not be
reported interchangeably.

\paragraph{Expected contrast with a targeted objective.}
Under a targeted loss every scored pixel is pushed toward the \emph{same}
class, so per-pixel gradients align and add. Untargeted losses have pixel
$A$ saying ``not road'' and pixel $B$ saying ``not sky'', and the two partly
cancel. That is the mechanistic reason the untargeted family converges on
one channel anyway --- and it also explains the sink-class variance of
Sec.~10.4.

% =====================================================================
\section{What to measure}

\subsection{Remote mIoU drop --- the headline}

\begin{equation}
  \Delta\mathrm{mIoU}_{\mathrm{rem}}
  = \mathrm{mIoU}\big(\mathcal{U}f(\mathbf{x}),\mathbf{y};\Omega\big)
  - \mathrm{mIoU}\big(\mathcal{U}f(\mathbf{x}^{\mathrm{adv}}),
     \mathbf{y};\Omega\big),
  \label{eq:drop}
\end{equation}
with the class mean taken over classes \emph{present in}
$\mathbf{y}|_\Omega$, so that both terms share a denominator
\citep{cordts2016cityscapes}. Counting a class merely because it is
\emph{predicted} makes the two means run over different class sets, and a
rare class entering or leaving the denominator shifts
$\Delta\mathrm{mIoU}$ by roughly $\mathrm{mIoU}/(n-1)$ independently of any
pixel-level change. This was an actual bug and it is worth one sentence in
the methods section.

\subsection{Flip rate}

\begin{equation}
  \phi = \frac{\big|\{(u,v)\in\Omega:\
     \arg\max_k \mathcal{U}f(\mathbf{x})_{k,uv}
     \neq \arg\max_k \mathcal{U}f(\mathbf{x}^{\mathrm{adv}})_{k,uv}\}\big|}
     {|\Omega|} .
  \label{eq:flip}
\end{equation}
$\phi$ is label-free, so it is the right instrument for the reach and
frequency probes, which must not reference any class.

\subsection{Influence beyond effect: margin and entropy by ring}

The winner margin at a pixel is
$\mathrm{top}_1-\mathrm{top}_2$ of the logits. Its change under attack shows
pixels the patch has pushed \emph{toward} a flip without completing one;
that region reliably extends past the flip zone, so the attack's
\emph{influence} is wider than its \emph{effect}. Normalised entropy
$H(\mathbf{p})/\ln C$ complements it: the margin measures the top-2 gap,
entropy measures uncertainty across all classes, and the two can disagree.
Both are reported over fixed absolute distance rings (50\,px steps to
1200\,px) so that ring $k$ is the same physical annulus in every image.

\subsection{Spectral efficiency --- the metric F1 requires}

Prediction change per unit of perceptual cost, per radial band $b$:
\begin{equation}
  \eta(b) = \frac{\phi(b)-\phi_0}{V_\beta(b)} ,
  \label{eq:efficiency}
\end{equation}
where $\phi(b)$ is the flip rate induced by band-limited noise confined to
$b$, $\phi_0$ is the unattacked baseline, and $V_\beta(b)$ is that
stimulus's visibility from Eq.~\eqref{eq:visibility}. \textbf{Quote $\eta$,
not $\phi$}: Sec.~10.1 is a statement about $\eta$ and is not supported as a
statement about $\phi$.

\subsection{Where the energy went}

Two radial summaries, and they answer different questions.
\begin{align}
  \mathrm{RAPSD}(b) &= \frac{1}{|A_b|}\sum_{f\in A_b}
     \big|\mathcal{F}\boldsymbol{\delta}(f)\big|^2 ,
     \label{eq:rapsd}\\
  \bar a(b) &= \frac{1}{|A_b|}\sum_{f\in A_b}
     \big|\mathcal{F}\boldsymbol{\delta}(f)\big| ,
     \qquad
  \mathrm{spend}(b) = \frac{\bar a(b)}{\bar B(b)} ,
  \label{eq:spend}
\end{align}
with $A_b$ the set of bins whose radial frequency falls in bin $b$.
Equation~\eqref{eq:rapsd} is \emph{power} \citep{ruzanski2011rapsd} and is
the right thing to compare against the natural-image power law
\citep{field1987}. Equation~\eqref{eq:spend} is \emph{amplitude}, because
the budget is an amplitude bound and the quantity worth plotting is the
ratio spent/allowed --- meaningful only if both sides are in the same units.

\paragraph{Bounded by CSF does not imply statistically inconspicuous.}
Natural spectra follow an approximate power law
\citep{field1987,falck2025fourier,dieleman2024spectral}, so energy placed
near Nyquist is \emph{atypical of the data} precisely where the eye is least
sensitive. Invisibility to a human and indistinguishability from natural
image statistics are distinct properties; only the former is constrained
here, and Eq.~\eqref{eq:rapsd} is what makes the latter measurable. A
detector trained on RAPSD is the obvious defence and should be named as
such.

\subsection{Uncertainty}

Report a mean with an interval and $n$, never a single run. Population runs
use a normal-approximation interval
$\bar x \pm 1.96\,\hat\sigma/\sqrt{n}$, chosen over a bootstrap
\citep{efron1993bootstrap} only because at $n$ in the hundreds they agree to
well inside the width that matters and an interval a reader can recompute
from $(\bar x,\hat\sigma,n)$ by hand is worth more than a slightly tighter
one. $n=1$ must return \emph{no} interval rather than a zero-width one: a
zero-width interval reads as ``measured exactly'', which is the opposite of
what one sample establishes. Pooled rates are pooled ratios --- numerator
and denominator summed over images, divided once --- with the per-image
spread reported alongside, never instead.

% =====================================================================
\section{Findings}
\label{sec:findings}

\subsection{F1 --- The network is spectrally flat; the gap is efficiency,
not effectiveness}

Measured on SegFormer MiT-B0, Cityscapes $512\times1024$, 5 images, patch
region, with band-limited noise injected at fixed RMS $=0.02$ (the
\emph{equal-amplitude} control, not the visibility-normalised run):
\begin{center}
\begin{tabular}{@{}lcc@{}}
\toprule
band ($\cpp$) & flip rate, remote & perceptual cost $V_\beta$ \\
\midrule
$0.00$--$0.02$ & $0.36\%$ & $70.7$ JND \\
$0.20$--$0.35$ & $0.53\%$ & $5.5$ JND \\
$0.35$--$0.50$ & $0.39\%$ & $1.5$ JND \\
\bottomrule
\end{tabular}
\end{center}
Every band moves the network about equally; their perceptual prices differ
by a factor of $\approx47$. Two corollaries, each a result in its own right:
\begin{itemize}
\item \textbf{The stride-4 patch embedding does not kill the high band.}
  That was the specific architectural risk this attack family was exposed
  to, and it did not materialise. This is the spectral counterpart to the
  ERF work: the ERF bounds reach \emph{spatially}, and no comparable
  spectral bound was found.
\item \textbf{The asymmetry lives entirely on the perceptual side.} The CSF
  constraint does not reduce what the attack can do; it reallocates a fixed
  perceptual budget to where it buys most. The original framing --- ``the
  network reads what the eye discards'' --- is half right: the network reads
  everything, and \emph{that} is what makes the CSF weighting useful.
\end{itemize}

\paragraph{What F1 does not support.}
Flip-rate differences between bands ($0.36$--$0.53\%$) sit inside their own
between-probe spread ($\pm0.32$--$0.63$), so they must not be reported as a
band ranking. Scene-dependence of the readable band is likewise
\emph{unsupported}: per-image peak bands came out $[3,4,2,4,3]$ from means
of $0.36$--$0.53\%$ against standard deviations of $0.32$--$0.63$, and the
tool now refuses that verdict when band separation is below the
between-probe spread. The hypothesis remains attractive --- it would explain
why a 150-epoch shared generator recovered $\approx0$ while per-image
optimisation reached $+16.31$ --- but it needs more probes and images before
it can be claimed. The visibility-normalised runs ($\tau\in\{0.25, 0.5, 1,
2\}$) were all \emph{inconclusive} against the probe's $0.5\%$ absolute
floor (maxima $0.036/0.108/0.241/0.425\%$); their band profile nevertheless
rose monotonically toward Nyquist in every $\tau$ and every scene,
$\approx40\times$ bottom to top. Report both the floor verdict and the
profile consistency.

\paragraph{Relation to the low-frequency attack literature.}
A body of work argues that adversarial perturbations restricted to
\emph{low} frequencies transfer better and survive input transformations
more robustly \citep{guo2020lowfreq,sharma2019spectral}. That is not in
tension with F1, and the distinction is worth stating explicitly because a
reader will otherwise assume one of the two must be wrong. Those results
optimise for \emph{transferability and robustness} under an $L_p$ budget,
where low frequencies are favoured because they survive resizing, JPEG and
model change. This work optimises for \emph{invisibility} under a
\emph{perceptual} budget, where low frequencies are the most expensive part
of the spectrum by a factor of $\approx47$ (F1) and the $L_p$ norm is not
the constraint at all. The two literatures agree that the spectrum is the
right place to think and disagree only about which end of it a given budget
makes attractive.

\subsection{F2 --- The budget was $\approx1.9\times$ too permissive}

Derived in Sec.~4.4 and measured over the full train split ($n=2975$). The
practical instruction: report \textbf{both} a nominal and a calibrated
$\tau$ column in every results table, so past runs are restatable by a
constant factor rather than invalidated. The enforced $\tau$ deliberately
keeps the legacy convention so the existing ladder stays comparable.

\subsection{F3 --- The constraint broke during optimisation, and
enforcement fixes it}

Sec.~5.5. Unenforced, realised visibility reached $2.77\times\tau$ (worst
seed $4.52\times$) on the legacy convention and $7.44\times$ (worst
$12.34\times$) on the locally measured one, at 1000 steps; at 400 steps it
was $1.57\times$ / $3.78\times$, i.e.\ the violation \emph{grows with run
length}. Enforced runs hold $0.2500$ exactly at every step.
\textbf{Instruction: \texttt{--csf\_enforce realised} for any number quoted
with a $\tau$ attached.}

\paragraph{The efficiency reading, which is the interesting part.}
Drop per unit of \emph{realised} visibility was better at 400 steps (best
$143.7$) than at 1000 (best $113.3$). The extra 600 steps buy drop by
spending perceptual budget, not by finding a cheaper attack. Any claim that
longer optimisation improves the attack must therefore be made per unit of
realised visibility, not per step.

\subsection{F4 --- Run length, and a bimodality that was not real}

Same configuration, five seeds per cell:
\begin{center}
\begin{tabular}{@{}lcc@{}}
\toprule
steps & $\Delta\mathrm{mIoU}_{\mathrm{rem}}$ & $\phi$ \\
\midrule
$400$  & $35.60\pm4.78$ & $63.1\pm6.7\%$ \\
$1000$ & $50.53\pm1.54$ & $98.95\pm0.33\%$ \\
$1000$, enforced & $45.95\pm3.57$ & --- \\
\bottomrule
\end{tabular}
\end{center}
At 400 steps the distribution looked bimodal (three runs $\approx39$, two
$\approx30$, nothing between). It was not two basins --- it was truncation:
every run was still descending when the step counter ran out. The
run-to-run instability that made two identical configurations look like a
code regression was mostly \emph{too few steps}; a flat learning rate (F6)
made it worse by never settling. Enforcement costs $4.58$ points.

\paragraph{Where the variance went.}
At 1000 steps the seeds agree on drop (relative sd $3.1\%$) and disagree
hard on cost (visibility sd $37.9\%$). Under enforcement the residual
variance is a \emph{sink-class choice}, not noise: a \texttt{building} sink
cannot flip building, which is $28.1\%$ of that frame, so its flip rate is
capped at $71.9\%$ (observed $68.66\%$, drop $40.86$), whereas a
\texttt{train} sink --- a class absent from the image --- flips building too
and reaches $99.4\%$ (the other four seeds: $47.22\pm2.48$). Forcing the
sink with a targeted objective is the way to remove that lottery, and it is
the cleanest available control for the implicit-class-selector finding (F7).

\subsection{F5 --- A fixed reference luminance costs almost nothing}

Making the CSF budget's reference luminance per-image buys only
$\approx13\%$ of budget at Nyquist. Over the full train split ($n=2975$),
p95/p5 budget ratio in sRGB code units:
\begin{center}
\begin{tabular}{@{}lccc@{}}
\toprule
 & low & mid & near-Nyquist \\
\midrule
train, centre (barten) & $1.63\times$ & $1.16\times$ & $1.13\times$ \\
train, road $(0.75,0.5)$ & $1.67\times$ & $1.16\times$ & $1.13\times$ \\
val, centre (barten) & $1.68\times$ & $1.17\times$ & $1.14\times$ \\
\bottomrule
\end{tabular}
\end{center}
Worst case over the \emph{full} range (min to max $Y$, a $93\times$ span in
linear luminance) is still only $1.39\times$ at Nyquist. Val tracks train
(median $Y$ $0.103$ vs $0.097$), so an $L_{\mathrm{ref}}$ fitted on train is
valid for val. Placement does not dominate: the road surface is $18\%$
darker in median than centre but only slightly wider in spread.

\paragraph{Mechanism --- and it is not ``Barten is flat''.}
Equation~\eqref{eq:cancellation}: the sRGB budget is proportional to
$c_{\mathrm{ref}}+0.055$, so gamma encoding has already performed most of
the luminance adaptation, the $2.4$ exponent cancelling against the same
exponent in the transfer slope. Barten's sensitivity then pushes back a
\emph{further} $0.59$--$0.85\times$, because a darker background is also a
less sensitive observer. Measuring in linear luminance instead reports
$7$--$13\times$ and is measuring in a space nothing in the attack lives in.
The SSO control confirms the decomposition independently: having no
luminance parameter, its ratio is flat at $1.930\times$ across all
frequencies, exactly $4.85\times$ (luminance) $\times\,0.398\times$ (sRGB
slope). Barten's frequency-dependent pushback is the entire difference
between $1.93$ and $1.13$.

\paragraph{Scope --- read before quoting this.}
This measures exactly \emph{one} term: the reference luminance the budget is
evaluated at. The per-image mode is content-adaptive in several other ways,
none of which this touches:
\begin{center}
\small
\begin{tabular}{@{}lccl@{}}
\toprule
mechanism & per-image today & measured by F5? \\
\midrule
base = the region the patch covers & yes & no \\
range fit, Eq.~\eqref{eq:fit} (per-image scale) & yes & no \\
the optimised residual itself & yes & no \\
placement (semantic / gradcam) & yes & no \\
CSF budget reference luminance & \textbf{no} (fixed $\mu=0.5$) & \textbf{yes}, and it buys $\approx13\%$ \\
\bottomrule
\end{tabular}
\end{center}
\textbf{Never write ``content-adaptivity has no perceptual justification''}
--- that over-reads this measurement by a wide margin. Argue for per-image
adaptivity on the residual, base, range fit and placement, and on attack
strength; do not argue for it on the grounds that a fixed reference
luminance mis-states visibility, because it barely does.

\subsection{F6 --- The parameterisation determines the learning rate}

Sec.~7.2. \texttt{csf} takes lr $0.2$, \texttt{universal\_csf} takes lr
$0.01$, and carrying one across to the other silently destroys the run in
one direction and freezes it in the other. Annealing is mandatory: a flat
learning rate swung $9.6$ mIoU points across four identical invocations.
Monitor $\pi$ (Eq.~\eqref{eq:atbound}) and $\mathrm{frac_{clip}}$
(Eq.~\eqref{eq:fracclip}) as first-class diagnostics, not as debug output.

\subsection{F7 --- The untargeted loss is an implicit class selector}

Sec.~8.5. $\approx97\%$ of remote flips on one channel, road~$\to$~car, with
car IoU down $\approx68$ points in every image. Report the flow table
alongside every mIoU number, and classify the result as \emph{structured}
(one flow $>50\%$) or \emph{diffuse}.

\subsection{F8 --- Reach is bounded spatially, and the bound is
architectural}

Prediction-change from a maximal random stimulus collapses below $1\%$ past
$\approx300$~px at $512\times1024$ and $\approx150$~px at $256\times512$;
$\approx50$--$80$k of $\approx450$k remote pixels are reachable. This is
label-free and target-independent by construction, which is what makes it a
valid control for the geometric factor rather than a correlate of the
attack. Restricting the loss support to the reachable set shrinks the
denominator $\approx5\times$; it cannot extend reach, and a null result
there is itself a test of the two-factor model.

\subsection{F9 --- Hard manifold constraints fail; soft and spectral ones
do not}

Constraining the patch to a 120-dimensional BigGAN latent --- a hard
manifold constraint --- found nothing adversarial inside it: $0.71\%$ flip
rate against $99.63\%$ for the unconstrained parameterisation on SegFormer,
across two architectures and both objectives. A 150-epoch shared conditional
generator likewise recovered $\approx0$ where per-image optimisation reached
$+16.31$. The spectral constraint of Sec.~5 is the opposite kind of object:
it keeps the full pixel parameterisation and reshapes rather than restricts
it, which is why it works where the latent constraint did not. This is the
argument for why the universal mode of Sec.~6 is the \emph{right} control
--- it isolates universality from parameterisation.

\subsection{The $\tau$ ladder, and where it stops meaning anything}

Below the knee of Eq.~\eqref{eq:fit}, realised visibility tracks $\tau$ and
the residual scales linearly:
\begin{center}
\begin{tabular}{@{}lcccc@{}}
\toprule
$\tau$ & $0.05$ & $0.10$ & $0.25$ & $0.50$ \\
residual RMS & $0.0049$ & $0.0098$ & $0.0245$ & $0.0490$ \\
\bottomrule
\end{tabular}
\end{center}
Above the knee the range fit binds, attained visibility saturates, and
$\tau$ ceases to control anything. The knee is reference-dependent --- it is
set by the headroom of $\mathbf{r}$ --- and lay near $\tau\approx0.4$ for
the references used here. \textbf{Usable range: $\tau\in[0.05,0.5]$.} Values
above the knee must not be reported as though achieved.

% =====================================================================
\section{Experimental design}
\label{sec:design}

\subsection{The reference configuration}

Everything below varies one axis against this fixed point. It is the
configuration the single-image result was promoted from, and it is chosen
rather than defaulted:
\begin{center}
\small
\begin{tabular}{@{}llL{7.3cm}@{}}
\toprule
setting & value & why this value \\
\midrule
architecture & SegFormer MiT-B0 & the model every finding above was measured
   on \citep{xie2021segformer} \\
resolution & $512\times1024$ & reach and band profiles are resolution-dependent
   (F8) \\
patch & $S=p=128$, $s=0.25$ & pinned; $S\neq p$ resamples the spectrum and
   invalidates $\tau$ (Sec.~6.4) \\
placement & \texttt{center} & keeps the footprint comparable across images;
   \texttt{gradcam} is a separate rung, not a default \\
mode & \texttt{csf --from\_image} & the base is the covered content
   (Sec.~5.6) \\
$\tau$ & $0.25$ & the rung every existing number is quoted at \\
enforcement & \texttt{--csf\_enforce realised} & mandatory for any number
   quoted with a $\tau$ (F3) \\
steps & $1000$ & $400$ truncates mid-descent and fakes bimodality (F4) \\
lr & $0.2$, cosine & correct for \texttt{csf} \emph{only}; flat swings
   $9.6$ points (F6) \\
objective & CosPGD & the main objective; CE and Tsallis are the axis \\
log interval & $25$ steps & 40 points for the convergence band, and the
   resolution at which ``best'' is checkpointed \\
\bottomrule
\end{tabular}
\end{center}

\subsection{The axes worth running, in order}

\begin{center}
\small
\begin{tabular}{@{}clL{7.6cm}@{}}
\toprule
\# & axis & what it settles \\
\midrule
1 & $\tau$ ladder $\{0.05,0.1,0.25,0.5\}$ &
    The deliverable is the \emph{curve}, not any operating point: strength
    against perceptibility, the role $\alpha$ plays in \citet{tan2021lap}.
    Stop at $0.5$; above the knee $\tau$ controls nothing. \\
2 & objective $\{$CE, CosPGD, Tsallis$\}$ &
    Sec.~\ref{sec:losses}. Hold $\tau$, placement, steps and lr fixed.
    Tsallis additionally sweeps $q\in\{0.5,1,2\}$ at fixed $\tau$. \\
3 & per-image vs universal &
    The adaptivity control (Sec.~6). Compare at matched \emph{nominal}
    $\tau$ \textbf{and} at matched \emph{realised} visibility --- the two
    modes enforce $\tau$ as an equality and an inequality respectively, so
    nominal matching alone is not a fair comparison. \\
4 & CSF-bounded vs unconstrained ceiling &
    A drop under a perceptual constraint is only interpretable next to the
    unconstrained ceiling \emph{on the same images}. \\
5 & CSF model $\{$barten, sso$\}$ &
    Shows the visibility claim does not depend on one CSF, and isolates the
    luminance decomposition (Sec.~3.3, F5). \\
6 & geometry sweep &
    Every visibility number is conditional on $\theta$. Report at least the
    desk-monitor default and the physical-patch geometry of
    Eq.~\eqref{eq:fromphysical}. \\
7 & architecture &
    F1 and F8 are single-architecture results. The frequency probe is
    label-free and dataset-independent, so it transfers to ADE20K weights
    and can bracket several backbones cheaply. \\
\bottomrule
\end{tabular}
\end{center}

\subsection{Sample size, pairing, and resumption}

\paragraph{One image is an anecdote.}
Neither a drop nor a realised-visibility number survives being quoted from a
single image, and neither survives being quoted from a single seed
(Sec.~10.4). Run the identical configuration over a population using the
\emph{same} function the single-image path uses, so that the two cannot
drift.

\paragraph{$n=100$, growing by permutation.}
$n=100$ supports a $\pm2$ point claim at the spread this attack has shown;
$n=500$ is the whole Cityscapes val set. The sample is a \emph{seeded
permutation truncated to $n$}, so $n=200$ is $n=100$ plus 100 new images and
a resumed run attacks only the new ones. Changing the seed instead throws
the first 100 away. Report the observed spread and the $n$ it would take to
support the claim being made, and raise $n$ only if that number says so.

\paragraph{Pair by image index.}
The constrained arm and the ceiling arm must use the same sample seed, and
the comparison must be the \emph{paired} difference. Given this attack's
scene-to-scene spread, the paired test is the only one that resolves a small
effect; an unpaired comparison of two means at $n=100$ will not.

\paragraph{Pilot before committing walltime.}
Five images first, so the walltime is arithmetic rather than a guess ---
1000 steps is $3.3\times$ the length the population loop was previously run
at.

\subsection{What every run must report}

A run that does not report these is not quotable.
\begin{enumerate}
\item $\Delta\mathrm{mIoU}_{\mathrm{rem}}$ with an interval and $n$
  (Eq.~\eqref{eq:drop}); never a single number.
\item Realised visibility $V^{\mathrm{real}}$ (Eq.~\eqref{eq:realvis}),
  \emph{both} nominal and calibrated (F2), as a distribution and not a mean.
\item The geometry $\theta$ in deg/px, and the Nyquist frequency it implies.
\item $\pi$ (Eq.~\eqref{eq:atbound}) for projected modes and
  $\mathrm{frac_{clip}}$ (Eq.~\eqref{eq:fracclip}) for clipped composites.
\item The convergence band: loss and drop against step, with an explicit
  verdict on whether the tail is flat. A run whose tail still climbs is
  truncated, not converged (F4).
\item The confusion flow table and the structured/diffuse verdict (F7).
\item The spend ratio by band, Eq.~\eqref{eq:spend}, against the budget it
  was allowed.
\item The seed count and the per-seed spread, for single-image results.
\end{enumerate}

\subsection{Claim discipline}

The measurements above license some sentences and forbid others. This table
is the one to keep next to the writing.
\begin{center}
\small
\begin{tabular}{@{}L{6.7cm}L{7.6cm}@{}}
\toprule
\textbf{Supported} & \textbf{Not supported} \\
\midrule
``Per unit of perceptual cost, near-Nyquist perturbation is
$\approx47\times$ more efficient than low-frequency perturbation'' (F1). &
``The network reads frequencies the eye cannot see.'' It reads all of them;
the difference is price. \\
``The readable band is a model property, consistent across the scenes
tested'' --- with the inconclusive verdict stated. &
``The readable band is scene-dependent.'' The separation is inside the
between-probe spread (F1). \\
``$\tau$ nominal $=0.25$; calibrated, that is $\approx0.47$'' (F2). &
Any $\tau$ quoted without saying which convention it is in. \\
``Under enforcement the constraint holds at $\tau$ exactly, at a cost of
$4.6$ mIoU points'' (F3). &
``The constraint holds by construction.'' It did not; it was violated
$2.8$--$7.4\times$ before enforcement. \\
``A fixed reference luminance costs $\approx13\%$ of budget at Nyquist''
(F5). &
``Content adaptivity has no perceptual justification.'' F5 measures one
term only. \\
``$\approx97\%$ of flips ran on one channel; the untargeted objective acts
as an implicit class selector'' (F7). &
``The untargeted attack causes broad-spectrum confusion.'' \\
``Below the modelled detection threshold of an average observer under the
stated geometry.'' &
``Imperceptible.'' That requires a study with human observers, which none of
this constitutes. \\
\bottomrule
\end{tabular}
\end{center}

\subsection{Sequencing, and the cheap results first}

Three of the most load-bearing measurements need no training and no GPU
time, and should be re-run whenever a constant changes: the footprint
luminance survey (F2, F5), the frequency-sensitivity probe (F1, label-free
so it transfers across weights), and the ERF probe (F8). Run those before
any ladder; a premise that fails there fails cheaply.

% =====================================================================
\section{Provenance, verification, and threats to validity}

\subsection{Provenance}

\begin{center}
\small
\begin{longtable}{@{}lL{6.6cm}L{5.0cm}@{}}
\toprule
\textbf{Sec.} & \textbf{Element} & \textbf{Source} \\
\midrule
\endhead
3.1 & cycles/pixel $\to$ cycles/degree & \citet[App.~D]{galinska2026csflow} \\
3.1 & screen and physical-patch geometry constructors & \novel{} \\
3.2 & Barten CSF; parameter values & \citet{barten2003}; params
      \citet[App.~A]{galinska2026csflow} \\
3.3 & Standard Spatial Observer, oblique effect &
      \citet{watson2005sso,watson2000five} \\
3.4 & sRGB transfer function and its slope & \citet{iec61966} \\
3.4 & luminance weights & \citet{itu709} \\
3.4 & pupil diameter, trolands & \citet[Ch.~2]{barten1999} \\
4.1 & Michelson contrast, $\kappa=2/\mu$ & \citet{michelson1927} \\
4.2 & \emph{inversion of CSF into an amplitude budget} & \novel{} \\
4.3 & \emph{patch-specific low-frequency cutoff} & \novel{} \\
4.4 & \emph{code-value calibration of $\kappa$ and the closed-form budget} &
      \novel{} \\
5.1 & \emph{smooth spectral reparameterisation} & \novel{} \\
5.2 & Minkowski pooling / probability summation &
      \citet{quick1974,robson1981,watson1979} \\
5.4 & \emph{uniform-rescale range fit} & \novel{} \\
5.5 & \emph{bisection enforcement on the composited residual} & \novel{} \\
5.6 & \emph{covered content as the residual base} & \novel{} \\
6.1 & universal perturbation framing & \citet{moosavi2017universal} \\
6.2 & \emph{envelope normalisation to pooled $\tau$} & \novel{} \\
6.5 & \emph{calibrated per-image visibility} & \novel{} \\
7.1 & localised patch application & \citet{brown2017patch} \\
7.1 & silhouette masking & \citet{tan2021lap} \\
7.2 & Adam; cosine annealing &
      \citet{kingma2015adam}; \citet{loshchilov2017sgdr} \\
8.1 & footprint exclusion & \citet{mirsky2023ipatch} \\
8.1 & effective receptive field & \citet{luo2016erf} \\
8.2 & untargeted CE attack lineage &
      \citet{goodfellow2015fgsm,madry2018pgd,xie2017dag} \\
8.2 & the ``effort drifts to won pixels'' pathology & \citet{gu2022segpgd} \\
8.3 & CosPGD objective; declared deviations &
      \citet{agnihotri2024cospgd}; deviations \novel{} \\
8.4 & Tsallis entropy & \citet{tsallis1988}; \citet{havrda1967} \\
8.4 & \emph{Tsallis entropy as an untargeted segmentation attack} & \novel{} \\
8.4 & entropy objectives, opposite sign &
      \citet{grandvalet2004,wang2021tent} \\
9.1 & mIoU protocol, Cityscapes & \citet{cordts2016cityscapes} \\
9.5 & RAPSD & \citet{ruzanski2011rapsd} \\
9.5 & natural-image spectral statistics &
      \citet{field1987,falck2025fourier,dieleman2024spectral} \\
9.6 & bootstrap, as the alternative not taken & \citet{efron1993bootstrap} \\
--- & LPIPS, evaluation only, never in the objective & \citet{zhang2018lpips} \\
--- & physical-world threat model, EOT, printability &
      \citet{athalye2018eot,sharif2016nps} \\
\bottomrule
\end{longtable}
\end{center}

\noindent
Twelve elements carry no external source and are contributions of this work.
From \citet{galinska2026csflow} this work takes a function, its constants, a
unit conversion and two scoping caveats --- not a method: their contribution
is a weighting over flow-matching timesteps, which has no analogue here.
Their $\alpha$, which interpolates
$w_{\mathrm{final}}=\alpha w_{\mathrm{CSFlow}}+(1-\alpha)w_{\mathrm{base}}$,
is \emph{not} $\tau$: it blends two weightings and bounds nothing.

\subsection{Bibliography status}

Entries were assembled from the verified bibliography of
\texttt{docs/csf\_patch\_procedure.tex} plus standard references, and carry
a DOI or an unambiguous venue wherever one exists. Three notes:
\begin{itemize}
\item \textbf{\texttt{galinska2026csflow}} --- the arXiv identifier is
  carried over unchanged from the companion document and has not been
  independently re-checked here. Verify it once before submission, since it
  is cited for the CSF parameters, the unit conversion, and two caveats.
\item \textbf{\texttt{gu2022segpgd}} --- SegPGD is ECCV \textbf{2022}. A
  comment in \texttt{validation.sh} attributes it to ECCV 2024; that is a
  typo worth fixing at the source so the wrong year does not propagate into
  the thesis.
\item \textbf{Tsallis as an attack objective} --- \citet{tsallis1988} and
  \citet{havrda1967} are cited for the \emph{entropy}. No source is claimed
  for its use as an adversarial segmentation objective; it is marked
  \novel{} in Sec.~12.1. If a prior use is found in the literature, the
  marking must change, and that is a search worth doing before the loss
  chapter is written.
\end{itemize}

\subsection{Threats to validity}

\begin{enumerate}
\item \textbf{$\tau$ is nominal, not a psychophysical JND.} Three
  approximations stand between them, each biasing it: the mean luminance in
  $\kappa$ is assumed rather than measured (quantified, F2); the Minkowski
  sum in Eq.~\eqref{eq:visibility} runs over DFT bins rather than perceptual
  channels; and Eq.~\eqref{eq:cpd} fixes a viewing geometry that
  \citet[Sec.~6]{galinska2026csflow} identify as needing calibration.
  $\tau$ should be reported as a relative axis.
\item \textbf{No human study.} Nothing here establishes imperceptibility.
  The most that Eq.~\eqref{eq:matched} supports is a consistency check by
  eye at the matched viewing distance.
\item \textbf{Single architecture, single dataset.} F1 and F8 are SegFormer
  results on Cityscapes. The probes that produce them are label-free and
  therefore cheap to replicate elsewhere; until they are, the claims are
  architecture-scoped.
\item \textbf{Digital threat model.} There is no expectation over
  transformations \citep{athalye2018eot} and no printability term
  \citep{sharif2016nps} in the objective. A physical patch would need both,
  and the geometry of Eq.~\eqref{eq:fromphysical} is the right starting
  point --- it is more favourable than the desk-monitor default, not less.
\item \textbf{Invisibility is not inconspicuousness.} The constraint bounds
  visibility to a modelled human observer. It does not bound detectability
  by a spectral classifier, and by placing energy near Nyquist it makes the
  patch \emph{atypical} of natural image statistics (Sec.~9.5). Name this
  defence rather than waiting for a reviewer to.
\item \textbf{The linearisation in Eq.~\eqref{eq:slope}} is first-order and
  holds for small residuals. At large $\tau$, trust the measured realised
  visibility rather than the budget.
\end{enumerate}

\bibliographystyle{plainnat}
\bibliography{frequency_optimised_findings}

\end{document}
